\chapter{Kuhn--Tucker Conditions and Nonlinear Programming}

Equality constraints assume the constraint always binds. In reality a budget
may be left unspent, a quota may not be reached, and consumption or input
levels must be non-negative. The Karush--Kuhn--Tucker (KKT) conditions generalise the method of Lagrange to
inequality constraints (the setting of nonlinear programming) and in particular they admit corner solutions
where some constraints are slack and the corresponding multiplier is zero.

Two ideas organise the chapter. \emph{Complementary slackness} couples each constraint to its multiplier:
at the optimum a constraint binds or its shadow price is zero. And for
\emph{convex programmes} (concave objective,
convex constraint functions) the KKT conditions are not merely necessary
but sufficient, and are tied to the \emph{saddle-point condition} of the Lagrangian. The solution routine is a
disciplined case analysis: guess which constraints bind, solve the
resulting equality-constrained system, and verify the remaining
inequalities and multiplier signs.

\section{Core Rules}

\begin{rulebox}{The nonlinear programming problem}
The standard form is
\[
  \max_{\mathbf{x}} \; f(\mathbf{x}) \quad \text{subject to} \quad
  g_i(\mathbf{x}) \leq b_i, \quad i = 1, \dots, m,
\]
usually together with non-negativity $x_j \geq 0$. Any smooth problem can
be brought to this form: an equality constraint is two inequalities
($g \leq b$ and $-g \leq -b$), non-negativity $x_j \geq 0$ corresponds to
$-x_j \leq 0$, and a minimisation is the maximisation of $-f$.
\end{rulebox}

\begin{rulebox}{The Kuhn--Tucker conditions (with complementary slackness)}
Form the Lagrangian
$\mathcal{L} = f(\mathbf{x}) + \sum_i \lambda_i\bigl(b_i - g_i(\mathbf{x})\bigr)$.
At a local maximum $\mathbf{x}^*$ (assuming a constraint qualification)
there exist multipliers $\lambda_i^* \geq 0$ such that:
\begin{itemize}[leftmargin=1.6em,itemsep=0.3ex]
  \item \textbf{Stationarity:} $\pd{\mathcal{L}}{x_j}(\mathbf{x}^*) = 0$
        for $j = 1, \dots, n$, i.e.\
        $\grad f(\mathbf{x}^*) = \sum_i \lambda_i^* \grad g_i(\mathbf{x}^*)$;
  \item \textbf{Primal feasibility:} $g_i(\mathbf{x}^*) \leq b_i$ for every
        $i$;
  \item \textbf{Dual feasibility:} $\lambda_i^* \geq 0$ for every $i$;
  \item \textbf{Complementary slackness:}
        $\lambda_i^*\,\bigl(b_i - g_i(\mathbf{x}^*)\bigr) = 0$ for every
        $i$. Equivalently: either the constraint binds ($g_i = b_i$) or its
        multiplier is zero.
\end{itemize}
Read-out: if constraint $i$ does not bind at the optimum, relaxing it
further has no effect on the value, so $\lambda_i^* = 0$. If it binds,
$\lambda_i^* \geq 0$ measures the shadow price of relaxing it, exactly as
in the equality case.
\end{rulebox}

\begin{rulebox}{Sign conventions for maximum problems with $g(\mathbf{x}) \leq b$}
Write every constraint as $g_i(\mathbf{x}) \leq b_i$ and form
\[
  \mathcal{L} = f(\mathbf{x})
  + \sum_i \lambda_i\bigl(b_i - g_i(\mathbf{x})\bigr),
  \qquad \lambda_i \geq 0 .
\]
Each $\lambda_i$ is non-negative, and if constraint $i$ binds it is the
shadow price of its right-hand side: $\partial V/\partial b_i = \lambda_i^*$.
Rules of thumb:
\begin{itemize}[leftmargin=1.6em,itemsep=0.3ex]
  \item a ``$\geq$'' constraint $g \geq b$ becomes $-g \leq -b$;
  \item for a \textbf{minimisation} $\min f$ s.t.\ $g \leq b$, the
        multipliers in $\mathcal{L} = f - \lambda(b - g)$ should also be
        $\geq 0$; note the sign. Many textbooks write all problems
        as maxima after flipping signs.
\end{itemize}
\end{rulebox}

\begin{rulebox}{Handling non-negativity constraints}
Two equivalent treatments of $x_j \geq 0$:
\begin{itemize}[leftmargin=1.6em,itemsep=0.3ex]
  \item \textbf{As explicit constraints:} write $-x_j \leq 0$ with
        multiplier $\mu_j \geq 0$; stationarity becomes
        $\pd{f}{x_j} - \sum_i \lambda_i \pd{g_i}{x_j} + \mu_j = 0$
        together with $\mu_j\,x_j = 0$.
  \item \textbf{Folded into stationarity:} drop the $\mu_j$ and replace the
        FOC by
        \[
          \pd{\mathcal{L}}{x_j} \leq 0, \qquad x_j \geq 0, \qquad
          x_j\,\pd{\mathcal{L}}{x_j} = 0 .
        \]
        In the interior ($x_j > 0$) this forces the usual
        $\pd{\mathcal{L}}{x_j} = 0$; at a corner ($x_j = 0$) only
        $\pd{\mathcal{L}}{x_j} \leq 0$ is required.
\end{itemize}
\end{rulebox}

\begin{rulebox}{When KKT is necessary: constraint qualification}
The KKT conditions are necessary at a local maximum only under a
\emph{constraint qualification}:
\begin{itemize}[leftmargin=1.6em,itemsep=0.3ex]
  \item \textbf{LICQ:} the gradients of the constraints \emph{active} at
        $\mathbf{x}^*$ are linearly independent;
  \item \textbf{Slater's condition} (for convex programmes): some feasible
        $\hat{\mathbf{x}}$ satisfies every inequality strictly,
        $g_i(\hat{\mathbf{x}}) < b_i$;
  \item affine (linear) constraints always qualify.
\end{itemize}
Standard economics problems (one budget plus non-negativity, or smooth
convex constraint sets) satisfy these. When the qualification fails, the
maximum may solve no KKT system: for $\max x$ s.t.\ $x^3 \leq 0$ the
optimum is $x^* = 0$, but $f'(0) = 1 \neq \lambda\,g'(0) = 0$ for every
$\lambda$.
\end{rulebox}

\begin{rulebox}{When KKT is sufficient: convex programmes}
The problem $\max f$ s.t.\ $g_i \leq b_i$ is a \emph{convex programme} if $f$ is concave and every constraint function
$g_i$ is convex, so that the feasible set is convex. Then:
\begin{itemize}[leftmargin=1.6em,itemsep=0.3ex]
  \item any point satisfying the KKT conditions is a \textbf{global
        maximum}: the conditions are sufficient, and no second-order
        (bordered-Hessian) check is needed;
  \item if $f$ is strictly concave, the maximiser is unique.
\end{itemize}
Linear objectives and linear constraints always qualify (linear functions
are both concave and convex). For a linear programme the Kuhn--Tucker
multipliers of the resource constraints are exactly its \textbf{dual
variables}, the shadow prices of the resources, and the KKT system is
primal feasibility, dual feasibility and complementary slackness side by
side: see the box ``LP duality and the Kuhn--Tucker conditions'' in
\chref{linprog}{7}. \textbf{Warning:} if the objective is convex
(Problem~4), maxima sit at vertices of the feasible set and stationary
points of the Lagrangian may be minima: KKT stationarity is then the wrong
tool.
\end{rulebox}

\begin{rulebox}{The saddle-point condition of the Lagrangian}
For $\max f(\mathbf{x})$ s.t.\ $g_i(\mathbf{x}) \leq b_i$ with
$\mathcal{L}(\mathbf{x}, \boldsymbol{\lambda})
= f(\mathbf{x}) + \sum_i \lambda_i\bigl(b_i - g_i(\mathbf{x})\bigr)$, a
pair $(\mathbf{x}^*, \boldsymbol{\lambda}^*)$ with
$\boldsymbol{\lambda}^* \geq \mathbf{0}$ is a \emph{saddle point} of the Lagrangian if
\[
  \mathcal{L}(\mathbf{x}, \boldsymbol{\lambda}^*)
  \;\leq\; \mathcal{L}(\mathbf{x}^*, \boldsymbol{\lambda}^*)
  \;\leq\; \mathcal{L}(\mathbf{x}^*, \boldsymbol{\lambda})
  \qquad \text{for all } \mathbf{x}
  \text{ and all } \boldsymbol{\lambda} \geq \mathbf{0},
\]
i.e.\ $\mathbf{x}^*$ maximises $\mathcal{L}(\cdot, \boldsymbol{\lambda}^*)$
and $\boldsymbol{\lambda}^*$ minimises
$\mathcal{L}(\mathbf{x}^*, \cdot)$ over
$\boldsymbol{\lambda} \geq \mathbf{0}$.
\begin{itemize}[leftmargin=1.6em,itemsep=0.3ex]
  \item If $(\mathbf{x}^*, \boldsymbol{\lambda}^*)$ is a saddle point, then
        $\mathbf{x}^*$ is feasible, complementary slackness holds, and
        $\mathbf{x}^*$ is a \textbf{global maximum}: no differentiability,
        concavity, or constraint qualification is needed (Problem~20).
  \item Conversely, in a convex programme satisfying Slater's condition,
        every solution $\mathbf{x}^*$ admits a multiplier vector
        $\boldsymbol{\lambda}^*$ such that
        $(\mathbf{x}^*, \boldsymbol{\lambda}^*)$ is a saddle point. For
        differentiable convex programmes, therefore: saddle point
        $\Leftrightarrow$ KKT point $\Leftrightarrow$ global maximum.
  \item If non-negativity $\mathbf{x} \geq \mathbf{0}$ is kept as a domain
        restriction rather than a constraint, the first inequality is only
        required for $\mathbf{x} \geq \mathbf{0}$.
\end{itemize}
\end{rulebox}

\begin{rulebox}{Solving strategy: systematic case analysis}
With $m$ inequality constraints there are at most $2^m$ binding/slack
patterns.
\begin{enumerate}[leftmargin=1.8em,itemsep=0.3ex]
  \item Guess which constraints bind (economic reasoning helps: a strictly
        increasing utility spends the whole budget).
  \item Impose the binding constraints as equalities, set the multipliers
        of the slack constraints to zero, and solve the reduced FOC system.
  \item Verify: (i) the omitted inequalities still hold at the candidate,
        (ii) the multipliers of the binding constraints are non-negative.
  \item If a check fails, switch the active set and re-solve.
\end{enumerate}
In a convex programme the first candidate that passes all checks is the
global maximum; otherwise compare $f$ across all surviving candidates.
\end{rulebox}

\begin{gallery}
\hspace*{\fill}%
\begin{tikzpicture}
  \begin{axis}[geoaxis, width=7.6cm, height=7.6cm,
      xmin=-2.6, xmax=4.6, ymin=-2.6, ymax=4.6,
      xtick={-2,1,3}, ytick={-2,1,4}]
    % feasible set: inside circle (x-1)^2+(y-1)^2<=9, x<=3, x+y>=1
    \addplot[draw=none, name path=upA, domain=-1.5616:2.5616, samples=201]
      {1+sqrt(9-(x-1)^2)};
    \addplot[draw=none, name path=linA, domain=-1.5616:2.5616, samples=2]{1-x};
    \region{upA}{linA}{}
    \addplot[draw=none, name path=upB, domain=2.5616:3, samples=101]
      {1+sqrt(9-(x-1)^2)};
    \addplot[draw=none, name path=loB, domain=2.5616:3, samples=101]
      {1-sqrt(9-(x-1)^2)};
    \region{upB}{loB}{}
    % boundaries
    \addplot[domain=0:360, samples=181]({1+3*cos(x)},{1+3*sin(x)});
    \addplot[domain=-2.4:4.4, samples=2]{1-x};
    \addplot[samples=2] coordinates {(3,-2.5)(3,4.5)};
    \addplot[only marks, mark=*, mark size=1.3pt, black, forget plot]
      coordinates {(3,1) (0,1) (1,4) (4,1)};
    \node[font=\tiny, anchor=south west] at (axis cs:3.05,1.08) {$P_{1}$};
    \node[font=\tiny, anchor=south east] at (axis cs:-0.06,1.08) {$P_{2}$};
    \node[font=\tiny, anchor=south west] at (axis cs:1.08,4.02) {$P_{3}$};
    \node[font=\tiny, anchor=south west] at (axis cs:4.06,1.08) {$P_{4}$};
  \end{axis}
\end{tikzpicture}%
\hspace*{\fill}
\gcap{The feasible set $C=\{(x-1)^{2}+(y-1)^{2}\le9,\ x\le3,\ x+y\ge1\}$,
with four test points.}
\end{gallery}

\begin{rulebox}{Four sign cases at a glance}
Write each constraint with its \emph{slack} $s_i\ge0$ when feasible
($s=b-g$ for $g\le b$, $s=g-b$ for $g\ge b$). For a maximum \emph{add}
$\lambda_is_i$, for a minimum \emph{subtract} it; then every multiplier is
$\ge0$ and complementary slackness reads $\lambda_is_i=0$.
\begin{center}\footnotesize
\begin{tabular}{llll}
  \hline
  problem & Lagrangian & stationarity & folded $x_j\ge0$\\
  \hline
  $\max f$, $g\le b$ & $f+\lambda(b-g)$ & $\grad f=\lambda\grad g$ & $\mathcal{L}_{x_j}\le0$\\
  $\max f$, $g\ge b$ & $f+\lambda(g-b)$ & $\grad f=-\lambda\grad g$ & $\mathcal{L}_{x_j}\le0$\\
  $\min f$, $g\ge b$ & $f+\lambda(b-g)$ & $\grad f=\lambda\grad g$ & $\mathcal{L}_{x_j}\ge0$\\
  $\min f$, $g\le b$ & $f+\lambda(g-b)$ & $\grad f=-\lambda\grad g$ & $\mathcal{L}_{x_j}\ge0$\\
  \hline
\end{tabular}
\end{center}
In the folded form, $x_j\ge0$ and $x_j\,\mathcal{L}_{x_j}=0$ complete the
conditions. Alternatively convert everything to ``$\max$ with $\le$''
($\min f=-\max(-f)$, $g\ge b\iff-g\le-b$) and use the first row only; both
routes give the same points and the same multipliers.
\end{rulebox}

\begin{rulebox}{Equality and inequality constraints together}
For $\max f$ s.t.\ $g_i\le b_i$ ($i\le m$) and $h_k=c_k$ ($k\le p$) use
\[
  \mathcal{L}=f+\sum_i\lambda_i(b_i-g_i)+\sum_k\mu_k(c_k-h_k),
  \qquad \lambda_i\ge0,\quad \lambda_i(b_i-g_i)=0,\quad
  \mu_k\ \text{free in sign},
\]
with stationarity $\grad f=\sum\lambda_i\grad g_i+\sum\mu_k\grad h_k$ and
\emph{no} complementary slackness for the equalities (they always bind).
Equivalent formulations of the same problem:
\begin{itemize}[leftmargin=1.6em,itemsep=0.3ex]
  \item an equality as two inequalities $h\le c$, $-h\le-c$ with
        multipliers $\mu^{+},\mu^{-}\ge0$: then $\mu=\mu^{+}-\mu^{-}$, and the
        pair is not unique (the two gradients are dependent);
  \item non-negativity folded into stationarity ($\mathcal{L}_{x_j}\le0$,
        $x_j\mathcal{L}_{x_j}=0$) or written as $-x_j\le0$ with multipliers
        $\nu_j\ge0$: then $\nu_j=-\mathcal{L}_{x_j}$ at the solution.
\end{itemize}
\end{rulebox}

\begin{rulebox}{Is the programme regular? A checking procedure}
A programme is \emph{regular} if a constraint qualification holds at
every feasible point; then every optimum is a KKT point. Check in this
order:
\begin{enumerate}[leftmargin=1.8em,itemsep=0.3ex]
  \item \textbf{Only linear constraints:} regular.
  \item \textbf{Convex constraint set} (each $g_i$ convex for $g_i\le b_i$):
        find one \textbf{Slater point} with every inequality strict; then
        regular.
  \item \textbf{Otherwise, LICQ pattern by pattern.} List every possible
        active set: each boundary curve on its own, and each corner where
        two (or more) boundaries meet. Compute the active gradients there
        and look for points where they vanish or are linearly dependent.
        If there are none, the programme is regular.
\end{enumerate}
\textbf{When it fails.} At a \emph{cusp} the active gradients are
dependent and the optimum may satisfy no KKT system. Example: on
$\{y\le(1-x)^{3},\ y\ge0\}$ the point $(1,0)$ maximises $x$, but there
$\grad\bigl(y-(1-x)^{3}\bigr)=(0,1)$ and $\grad(-y)=(0,-1)$ are parallel,
and $(1,0)=\lambda_1(0,1)+\lambda_2(0,-1)$ is impossible.

\textbf{Uniqueness of multipliers.} If the active gradients are linearly
independent, the stationarity equation
$\grad f=\sum\lambda_i\grad g_i$ has at most one solution: the multipliers
are unique. If they are dependent (for instance three active constraints at
a corner in the plane, or an equality written as two inequalities), the
multipliers form a whole family, typically a one-parameter set, and the
shadow-price reading of each single $\lambda_i$ is lost.
\end{rulebox}

\begin{rulebox}{Convex programmes in minimisation form}
$\min f$ s.t.\ $g_i\le b_i$, $h_k\ge c_k$, $e_l=d_l$ is a \emph{convex
programme} if $f$ is \textbf{convex}, each $g_i$ (in a $\le$ constraint) is
\textbf{convex}, each $h_k$ (in a $\ge$ constraint) is \textbf{concave},
and each $e_l$ is \textbf{affine}. Then the feasible set is convex and every
KKT point is a global minimum. For a maximum, $f$ must be concave instead,
with the same conditions on the constraints. Consequences:
\begin{itemize}[leftmargin=1.6em,itemsep=0.3ex]
  \item non-convexity can come from the \emph{feasible set} (outside a
        circle: $x^{2}+y^{2}\ge1$ has a convex, not concave, left side; a
        hyperbola region $x^{2}-y^{2}\ge1$; the region under a convex
        parabola $y\le x^{2}$) or from the \emph{objective};
  \item the same feasible set can give a convex programme for $\min f$ and
        a non-convex one for $\max f$: minimising $x^{2}+y^{2}$ on a disc
        is convex, maximising it is not;
  \item to show convexity of a non-quadratic objective, check that its
        Hessian is PSD everywhere (\chref{optimization}{3}).
\end{itemize}
In a non-convex programme KKT points are only candidates: some may be
saddle points, local optima or even the opposite extremum, so compare the
objective at all of them (and at points where regularity fails).
\end{rulebox}

\begin{rulebox}{Graphical solution of a nonlinear programme}
In two variables:
\begin{enumerate}[leftmargin=1.8em,itemsep=0.3ex]
  \item Draw the feasible set: each boundary curve, and shade where all
        constraints hold.
  \item Draw level curves of $f$: circles around the unconstrained optimum
        for a distance objective, parallel lines for a linear one,
        parabolas or hyperbolas otherwise.
  \item Move the level in the improving direction until the last contact
        with the feasible set. The contact is a \emph{tangency} with one
        boundary curve or a \emph{corner} where two meet.
  \item Read off which constraints bind there, set the others' multipliers
        to zero, and solve the reduced KKT system for the point and the
        multipliers; check their signs.
\end{enumerate}
Do it twice, for the minimum and for the maximum: on a compact feasible set
both exist (Weierstrass), and they usually sit on different parts of the
boundary. The figure below shows both for a distance objective on a quarter
disc.
\end{rulebox}

\begin{gallery}
\hspace*{\fill}%
\begin{tikzpicture}
  \begin{axis}[contouraxis, width=6.4cm, height=6.4cm,
      xmin=-0.6, xmax=4.4, ymin=-0.6, ymax=4.4, xtick={1,2,3}, ytick={1,2,3}]
    \addplot[draw=none, name path=arc, domain=0:2.8284, samples=101]{sqrt(8-x^2)};
    \addplot[draw=none, name path=ax, domain=0:2.8284, samples=2]{0};
    \region{arc}{ax}{}
    \addplot[black, thick, domain=0:2.8284, samples=101]{sqrt(8-x^2)};
    \addplot[black, thick, samples=2] coordinates {(0,0)(2.8284,0)};
    \addplot[black, thick, samples=2] coordinates {(0,0)(0,2.8284)};
    \foreach \r in {1.4142, 2.8284, 4.2426} {
      \addplot[dashed, domain=0:360, samples=121]({3+\r*cos(x)},{3+\r*sin(x)}); }
    \addplot[only marks, mark=*, mark size=1.4pt, black, forget plot]
      coordinates {(2,2) (0,0) (3,3)};
    \node[font=\tiny, anchor=south west] at (axis cs:2.02,2.02) {min};
    \node[font=\tiny, anchor=north west] at (axis cs:0.03,-0.03) {max};
    \node[font=\tiny, anchor=south west] at (axis cs:3.02,3.02) {$(3,3)$};
  \end{axis}
\end{tikzpicture}%
\hspace*{\fill}
\gcap{$f=(x-3)^{2}+(y-3)^{2}$ on the quarter disc
$\{x^{2}+y^{2}\le8,\ x\ge0,\ y\ge0\}$: level circles around $(3,3)$ of
radius $\sqrt2$, $2\sqrt2$, $3\sqrt2$. The smallest touching circle gives
the minimum $2$ at the tangency point $(2,2)$ (only the arc binds); the
largest one gives the maximum $18$ at the corner $(0,0)$ (both axes bind).}
\end{gallery}

\begin{rulebox}{Testing whether a given point can be optimal}
\begin{enumerate}[leftmargin=1.8em,itemsep=0.3ex]
  \item \textbf{Feasibility:} check every constraint. An infeasible point
        is out.
  \item \textbf{Active set:} note which constraints hold with equality.
        Set the multipliers of all slack constraints to zero.
  \item \textbf{Stationarity:} solve for the active multipliers. No
        solution: not a KKT point (if the programme is regular there, not
        optimal either).
  \item \textbf{Signs:} a negative multiplier rules the point out; all
        $\ge0$ makes it a KKT point.
  \item \textbf{Optimal?} In a convex programme a KKT point is a global
        optimum. Otherwise it is only a candidate: compare with the other
        candidates.
\end{enumerate}
\end{rulebox}

\begin{rulebox}{Standard forms and variable substitutions}
Besides the constraint rewrites above, variables can be brought to sign
form, as for linear programmes (\chref{linprog}{3}):
\begin{itemize}[leftmargin=1.6em,itemsep=0.3ex]
  \item a \textbf{non-positive} variable $z\le0$: substitute $z=-z'$,
        $z'\ge0$;
  \item a \textbf{free} variable $x$: substitute $x=x^{+}-x^{-}$ with
        $x^{+},x^{-}\ge0$, everywhere it occurs, also inside a nonlinear
        objective or constraint.
\end{itemize}
Two common target forms: ``$\max F$ s.t.\ $G_i(\mathbf{x})\ge0$,
$\mathbf{x}\ge\mathbf{0}$'' and ``$\min F$ s.t.\ $G_i(\mathbf{x})\le0$,
$\mathbf{x}\ge\mathbf{0}$''. An equality becomes two inequalities.
\end{rulebox}

\begin{rulebox}{Lagrange versus Kuhn--Tucker; existence}
On a feasible set with bounds, the Lagrange method applied to the equality
part alone finds only tangency points: it can deliver the minimum and miss
a maximum that sits at a bound (a corner), where the KKT conditions hold
with positive bound multipliers. A complete solution therefore checks the
corners as well. On an \emph{unbounded} feasible set an optimum need not
exist. If the programme is regular, an optimum must be a KKT point, so
\emph{showing that no point satisfies the KKT system proves that no optimum
exists}.
\end{rulebox}

\section{Worked Example}

\begin{worked}{Maximise $f(x,y)=xy$ subject to $x+y\le10$, $x\le4$, $x\ge0$, $y\ge0$}
\wstep{Step 1. Write the Lagrangian with one multiplier per inequality.}
\[
  \mathcal L=xy+\lambda(10-x-y)+\mu(4-x) .
\]
The Kuhn--Tucker conditions for a maximum with $\le$ constraints are
\[
  \mathcal L_{x}=y-\lambda-\mu\le0,
  \qquad
  \mathcal L_{y}=x-\lambda\le0,
\]
with equality wherever the corresponding variable is positive, together with
\[
  \lambda\ge0,\quad \mu\ge0,
  \qquad
  \lambda(10-x-y)=0,
  \qquad
  \mu(4-x)=0 .
\]
The last two are \textbf{complementary slackness}: a constraint that is not
binding has a zero multiplier, and a positive multiplier forces its
constraint to bind.
\wtrap{The signs matter here. For a maximum with $g(\mathbf x)\le b$ the
multipliers are $\ge0$; flip to $\ge$ constraints or to minimisation and
they flip too. Fix one convention --- ``maximise, all constraints written as
$\le$, multipliers $\ge0$'' --- and convert every problem into it rather
than memorising four sign tables.}

\wstep{Step 2. Enumerate the cases systematically.}
With two inequality constraints there are four combinations of
binding/slack. Work through them in order; do not guess.

\wstep{Step 3. Case A --- both slack ($\lambda=\mu=0$).}
Then $\mathcal L_{y}=x=0$ and $\mathcal L_{x}=y=0$, giving $(0,0)$ with
$f=0$. Feasible, and a legitimate KKT point --- but it is the \emph{minimum}
of $f$ on this set, so keep it as a candidate and continue.

\wstep{Step 4. Case B --- only the budget binds ($\mu=0$, $\lambda>0$).}
Then $x=\lambda$ and $y=\lambda$, so $x=y$; with $x+y=10$ this gives
$x=y=5$. But $x=5$ violates $x\le4$: \textbf{infeasible}. Discard.
\wtrap{A case that produces an infeasible point is discarded, not
``adjusted''. The temptation to push $x$ back to $4$ and keep $y=5$ is
exactly the error --- that point does not satisfy the equations of this
case, and it belongs to a different case.}

\wstep{Step 5. Case C --- only $x\le4$ binds ($\lambda=0$, $\mu>0$).}
Then $\mathcal L_{y}=x-0=0$ forces $x=0$, contradicting $x=4$. Discard.

\wstep{Step 6. Case D --- both bind.}
$x=4$ and $x+y=10$ give $y=6$. Now solve for the multipliers:
\[
  \mathcal L_{y}=x-\lambda=0\Rightarrow\lambda=4 ;
  \qquad
  \mathcal L_{x}=y-\lambda-\mu=0\Rightarrow 6-4-\mu=0\Rightarrow\mu=2 .
\]
Both multipliers are $\ge0$ as required. \checkmark\ Feasible, all
conditions satisfied:
\[
  \boxed{(x^{*},y^{*})=(4,6),\qquad f^{*}=24 .}
\]

\wstep{Step 7. Check the multipliers' signs --- this is the test, not a formality.}
Had $\mu$ come out negative, Case~D would be rejected: a negative multiplier
on a $\le$ constraint in a maximisation problem means the constraint is
pushing the wrong way, i.e.\ removing it would \emph{lower} the objective,
which is impossible. The sign check is what eliminates spurious cases.

\wstep{Step 8. Confirm it is the global maximum.}
Compare the surviving candidates: $(0,0)$ with $f=0$ and $(4,6)$ with
$f=24$. The feasible set is closed and bounded (a quadrilateral with
vertices $(0,0),(4,0),(4,6),(0,10)$), so by Weierstrass a maximum exists and
must be among the KKT points. Hence $(4,6)$ is global.
\wtrap{KKT conditions are \emph{necessary} under a constraint qualification,
and sufficient only for a concave objective over a convex set. Here $xy$ is
\emph{not} concave --- its Hessian is indefinite --- so sufficiency does not
apply and the comparison in this step is doing real work. Skipping it would
leave $(0,0)$ and $(4,6)$ equally certified.}

\wstep{Step 9. Read the multipliers.}
$\lambda^{*}=4$: relaxing the budget from $10$ to $11$ raises $f$ by about
$4$. $\mu^{*}=2$: relaxing $x\le4$ to $x\le5$ raises $f$ by about $2$.
Check the first directly --- with $x+y\le11$ and $x\le4$ the optimum is
$(4,7)$ with $f=28$, exactly $4$ more. \checkmark\ Shadow prices work in the
inequality-constrained case exactly as they did in \chref{optimization}{8}.

\wstep{Step 10. The recipe.}
\begin{center}\footnotesize
\begin{tabular}{ll}
  \hline
  1 & write $\mathcal L$; state the sign convention\\
  2 & list stationarity, feasibility, sign, complementary slackness\\
  3 & enumerate all $2^{k}$ binding/slack cases for $k$ constraints\\
  4 & in each: solve, then test feasibility \emph{and} multiplier signs\\
  5 & compare the survivors; invoke concavity or Weierstrass to conclude\\
  \hline
\end{tabular}
\end{center}
With three or four constraints the enumeration grows, but economic reasoning
usually rules out most cases at a glance --- a budget constraint with
monotone preferences always binds, and a non-negativity constraint on a good
with infinite marginal utility at zero never does.
\end{worked}

\section{Practice Problems}

Problems 1--15 drill the KKT system and the case analysis of binding and
slack constraints; Problems 16--20 add convex-programme sufficiency and the
saddle-point condition.

\begin{problems}
  \item Maximise $f(x, y) = x y$ on $x + y \leq 8$, $x \geq 0$, $y \geq 0$.
  \item Maximise $f(x, y) = 4 x + 3 y$ on $x + y \leq 6$,
        $2 x + y \leq 8$, $x, y \geq 0$.
  \item Maximise $U(x, y) = x + \ln(1 + y)$ on $x + y \leq 10$,
        $x, y \geq 0$. \emph{Test:} is the budget tight?
  \item Maximise $f(x, y) = x^2 + y^2$ on $x + y \leq 1$, $x, y \geq 0$.
        \emph{Note:} the objective is convex; maxima are at vertices.
  \item Minimise $f(x, y) = (x - 3)^2 + (y - 4)^2$ on $x + y \leq 5$,
        $x, y \geq 0$.
  \item Maximise $U(x, y) = \ln(x) + \ln(y)$ on $x + y \leq 10$,
        $x \geq 1$, $y \geq 1$. Is the lower bound on $x$ binding? On $y$?
  \item Maximise $f(x, y) = 2 x + y$ on $x^2 + y^2 \leq 5$.
  \item A monopolist chooses $q \geq 0$ to maximise
        $\pi(q) = (10 - q) q - 4 q = 6 q - q^2$. (Use the non-negativity
        constraint explicitly.) Find $q^*$.
  \item Maximise $U = x y$ on $p_x x + p_y y \leq M$, $x \geq 0$,
        $y \geq 0$ (general parameters $p_x, p_y, M > 0$). Show the budget
        binds.
  \item Maximise $f(x, y) = x + 2 y$ on $x^2 + y^2 \leq 4$ and $y \leq 1$.
  \item Minimise $f(x) = x^2 - 4 x + 7$ on $x \geq 5$. (Note: the
        unconstrained optimum is at $x = 2$, infeasible.)
  \item Maximise $U(x_1, x_2) = x_1 + 2 \ln(1 + x_2)$ on
        $x_1 + x_2 \leq 5$, $x_1, x_2 \geq 0$.
  \item Maximise $f(x, y) = \ln x + \ln y$ on $x + y \leq 10$,
        $x \geq 1$, $y \geq 4$. Which constraints bind?
  \item Maximise $U = x^{1/2} + y^{1/2}$ on $x + 2 y \leq 12$,
        $x \geq 0$, $y \geq 0$.
  \item Maximise $f(x, y) = 6 x + 4 y - x^2 - y^2$ on
        $x \geq 0$, $y \geq 0$ (no upper-bound constraint).
  \item Show that
        \[
          \max\; f(x, y) = -(x - 2)^2 - (y - 1)^2
          \quad \text{s.t.} \quad x + 2 y \leq 2, \; x \geq 0, \; y \geq 0
        \]
        is a convex programme. Solve it via the KKT conditions and explain
        why no second-order check is needed.
  \item Maximise $U(x, y) = 4 \ln(1 + x) + y$ on $x + y \leq 2$,
        $x, y \geq 0$. Carry out a systematic case analysis (budget
        binding/slack, each variable zero or positive) and verify that the
        problem is a convex programme, so the KKT point you find is the
        global maximum.
  \item A firm's profit from two products is
        $\pi(x, y) = 12 x - x^2 + 8 y - y^2$; capacity restricts total
        output to $x + y \leq 8$, with $x, y \geq 0$. (a)~Solve by KKT and
        use convexity to certify a global maximum. (b)~Repeat with capacity
        $12$ instead of $8$ and interpret the multiplier via complementary
        slackness.
  \item For $\max f(x) = -x^2 + 6 x$ s.t.\ $x \leq 2$: find the KKT point
        $(x^*, \lambda^*)$ and verify \emph{directly} that it is a saddle
        point of the Lagrangian, i.e.\
        $\mathcal{L}(x, \lambda^*) \leq \mathcal{L}(x^*, \lambda^*)
        \leq \mathcal{L}(x^*, \lambda)$ for all $x \in \R$ and all
        $\lambda \geq 0$.
  \item (Saddle-point theorem.) Let
        $\mathcal{L}(\mathbf{x}, \lambda)
        = f(\mathbf{x}) + \lambda\bigl(b - g(\mathbf{x})\bigr)$ and suppose
        the pair $(\mathbf{x}^*, \lambda^*)$ with $\lambda^* \geq 0$
        satisfies the saddle-point condition
        $\mathcal{L}(\mathbf{x}, \lambda^*)
        \leq \mathcal{L}(\mathbf{x}^*, \lambda^*)
        \leq \mathcal{L}(\mathbf{x}^*, \lambda)$ for all $\mathbf{x}$ and
        all $\lambda \geq 0$. Prove: (i)~$\mathbf{x}^*$ is feasible,
        (ii)~$\lambda^*\bigl(b - g(\mathbf{x}^*)\bigr) = 0$, and
        (iii)~$\mathbf{x}^*$ solves $\max f$ s.t.\ $g \leq b$. (No
        differentiability or convexity assumptions are needed.)
\end{problems}

\section{Solutions}

\begin{solutions}
  \item Guess an interior solution with the budget binding.
        $\mathcal{L} = x y + \lambda(8 - x - y)$. FOCs:
        \[
          \pd{\mathcal{L}}{x} = y - \lambda = 0, \qquad
          \pd{\mathcal{L}}{y} = x - \lambda = 0
          \;\Rightarrow\; x = y .
        \]
        Sum in the budget: $2 x = 8 \Rightarrow x^* = y^* = 4$, with
        $\lambda^* = 4 \geq 0$ \checkmark. Non-negativity is slack
        ($\mu_x = \mu_y = 0$). $f^* = 16$.

  \item Linear programming. The vertices of the feasible polygon are
        $(0, 0)$, $(4, 0)$, $(2, 4)$, $(0, 6)$. Evaluating $f$ there gives
        $0$, $16$, $20$, $18$: the optimum is at $(2, 4)$ with $f^* = 20$.
        Both inequality constraints bind; non-negativity is slack at this
        vertex. Multipliers from stationarity:
        \[
          \pd{\mathcal{L}}{x} = 4 - \lambda_1 - 2 \lambda_2 = 0, \qquad
          \pd{\mathcal{L}}{y} = 3 - \lambda_1 - \lambda_2 = 0 .
        \]
        Subtracting: $1 = \lambda_2$, hence $\lambda_1 = 2$. Both $\geq 0$.
        \checkmark

  \item Guess the budget binds with $x, y > 0$.
        $\mathcal{L} = x + \ln(1 + y) + \lambda(10 - x - y)$. FOCs:
        $1 = \lambda$ and $1/(1 + y) = \lambda = 1 \Rightarrow y^* = 0$,
        but we assumed $y > 0$. Switch the case: try $y = 0$ binding. Then
        $x = 10$, $y = 0$, $f^* = 10$. Multiplier for the budget:
        $\lambda = 1$ (from the FOC on $x$). Multiplier for $y \geq 0$:
        $\mu_y = \lambda - 1/(1 + 0) = 1 - 1 = 0$. So the corner $y = 0$
        with $\mu_y = 0$ is consistent (a boundary solution with vanishing
        multiplier, non-strict KKT). $x^* = 10$, $y^* = 0$, $U^* = 10$.
        The budget is tight.

  \item A convex objective on a compact convex set: the maximum is at a
        vertex of the simplex, $(0, 0)$, $(1, 0)$, $(0, 1)$. The
        $f$-values are $0$, $1$, $1$. \textbf{Both} $(1, 0)$ and $(0, 1)$
        are optima with $f^* = 1$. (A Lagrange-style FOC would only flag
        $(0, 0)$ as a critical point of $f$, which is a \emph{minimum}
        here.)

  \item The unconstrained minimum of the paraboloid is at $(3, 4)$, with
        $3 + 4 = 7 > 5$: infeasible. So the budget binds. Use the form
        $\mathcal{L} = (x - 3)^2 + (y - 4)^2 + \mu(x + y - 5)$ with
        $\mu \geq 0$ (the KKT-standard form for $g(x,y) - b \leq 0$ in a
        minimisation). FOCs:
        \[
          2(x - 3) + \mu = 0, \qquad 2(y - 4) + \mu = 0
          \;\Rightarrow\; 2(x - 3) = 2(y - 4)
          \;\Rightarrow\; y = x + 1 .
        \]
        Constraint: $x + y = 5 \Rightarrow 2 x + 1 = 5 \Rightarrow
        x^* = 2$, $y^* = 3$. Both $\geq 0$ \checkmark.
        $\mu^* = -2(x - 3) = -2(-1) = 2 \geq 0$ \checkmark.
        $f^* = 1 + 1 = 2$.

  \item Guess an interior solution with $x, y > 1$ and the budget binding.
        $\mathcal{L} = \ln x + \ln y + \lambda(10 - x - y)$. FOCs:
        $1/x = \lambda$, $1/y = \lambda \Rightarrow x = y = 5$. Both
        $> 1$: the lower bounds are slack with multipliers zero.
        $\lambda^* = 1/5$. $U^* = 2 \ln 5$.

  \item Guess the circle constraint binding.
        $\mathcal{L} = 2 x + y + \lambda(5 - x^2 - y^2)$. FOCs:
        \[
          2 = 2 \lambda x, \qquad 1 = 2 \lambda y
          \;\Rightarrow\; x = 1/\lambda, \quad y = 1/(2 \lambda) .
        \]
        Constraint: $1/\lambda^2 + 1/(4 \lambda^2) = 5/(4 \lambda^2) = 5
        \Rightarrow \lambda^2 = 1/4 \Rightarrow \lambda = 1/2$ (positive
        for a max). $x^* = 2$, $y^* = 1$. $f^* = 5$.

  \item The constraint is $-q \leq 0$, i.e.\ $q \geq 0$.
        $\mathcal{L} = 6 q - q^2 + \mu q$ with $\mu \geq 0$ and
        complementary slackness $\mu q = 0$. FOC: $6 - 2 q + \mu = 0$. Try
        $\mu = 0$: $q^* = 3 > 0$, so $\mu^* = 0$ confirms the interior
        solution. $\pi^* = 9$.

  \item Lagrangian
        $\mathcal{L} = x y + \lambda(M - p_x x - p_y y)
        + \mu_x x + \mu_y y$. FOCs: $y = \lambda p_x - \mu_x$,
        $x = \lambda p_y - \mu_y$. Since the utility $x y$ is strictly
        increasing in both arguments (in the interior), the consumer always
        wants to spend everything: the budget binds ($\lambda > 0$). And
        $x = 0$ would give $U = 0$; one can always do better by buying some
        of both. So $\mu_x = \mu_y = 0$. Then $y = \lambda p_x$,
        $x = \lambda p_y$. Substituting into the budget:
        \[
          p_x \lambda p_y + p_y \lambda p_x = 2 p_x p_y \lambda = M
          \;\Rightarrow\; \lambda^* = \frac{M}{2 p_x p_y} .
        \]
        Hence $x^* = M/(2 p_x)$, $y^* = M/(2 p_y)$.

  \item Two constraints: $g_1\colon x^2 + y^2 \leq 4$ and
        $g_2\colon y \leq 1$. The unconstrained maximum of $x + 2 y$ is
        unbounded, so $g_1$ must bind.

        \emph{Case A ($g_2$ binds too):} $y = 1$,
        $x^2 = 3 \Rightarrow x = \sqrt{3}$ (the positive root maximises).
        $f = \sqrt{3} + 2 \approx 3.73$. Multipliers:
        $\mathcal{L} = x + 2 y + \lambda_1(4 - x^2 - y^2)
        + \lambda_2(1 - y)$. FOCs: $1 = 2 \lambda_1 x \Rightarrow
        \lambda_1 = 1/(2 \sqrt{3})$;
        $2 = 2 \lambda_1 y + \lambda_2 = 2 \lambda_1 + \lambda_2
        \Rightarrow \lambda_2 = 2 - 1/\sqrt{3} \approx 1.42$. Both
        $\geq 0$ \checkmark.

        \emph{Case B ($g_2$ slack, $y < 1$):} FOCs
        $1 = 2 \lambda_1 x$, $2 = 2 \lambda_1 y \Rightarrow y = 2 x$. On
        the circle: $x^2 + 4 x^2 = 4 \Rightarrow x = 2/\sqrt{5}$,
        $y = 4/\sqrt{5} \approx 1.79 > 1$, infeasible.

        So the binding-both case is correct: $x^* = \sqrt{3}$, $y^* = 1$,
        $f^* = 2 + \sqrt{3}$.

  \item Minimisation on $x \geq 5$. The unconstrained minimum at $x = 2$
        is infeasible. The constraint is $-x + 5 \leq 0$, i.e.\
        $5 - x \leq 0$, so the minimum sits on the boundary: $x^* = 5$.
        $f(5) = 25 - 20 + 7 = 12$. KKT: stationarity $f'(x) = \mu$ where
        $\mu \geq 0$ is the multiplier on $5 - x \leq 0$ (for a min).
        $f'(5) = 6 = \mu \geq 0$. \checkmark

  \item $\mathcal{L} = x_1 + 2 \ln(1 + x_2) + \lambda(5 - x_1 - x_2)
        + \mu_1 x_1 + \mu_2 x_2$. Try an interior solution with the budget
        binding: FOCs $1 = \lambda$ and
        $2/(1 + x_2) = \lambda = 1 \Rightarrow x_2 = 1$. Then $x_1 = 4$.
        Both $\geq 0$ \checkmark, so $\mu_1 = \mu_2 = 0$.
        $U^* = 4 + 2 \ln 2$.

  \item Try the interior (budget binding, both lower bounds slack): FOC
        $1/x = 1/y \Rightarrow x = y = 5$. But $y = 5 > 4$ \checkmark and
        $x = 5 > 1$ \checkmark, so both lower bounds are slack with
        multipliers zero. $\lambda^* = 1/5$. $U^* = 2 \ln 5$. Confirmed by
        checking: at $(5, 5)$ no constraint other than the budget is
        binding.

  \item The budget should bind (utility strictly increasing).
        $\mathcal{L} = x^{1/2} + y^{1/2} + \lambda(12 - x - 2 y)$. FOCs:
        \[
          \frac{1}{2 \sqrt{x}} = \lambda, \qquad
          \frac{1}{2 \sqrt{y}} = 2 \lambda = \frac{1}{\sqrt{x}}
          \;\Rightarrow\; \sqrt{x} = 2 \sqrt{y}
          \;\Rightarrow\; x = 4 y .
        \]
        Constraint: $4 y + 2 y = 6 y = 12 \Rightarrow y^* = 2$,
        $x^* = 8$. $\lambda^* = 1/(2 \sqrt{8}) = 1/(4 \sqrt{2})$. Both
        variables positive, lower bounds slack.
        $U^* = \sqrt{8} + \sqrt{2} = 2 \sqrt{2} + \sqrt{2} = 3 \sqrt{2}$.

  \item Unconstrained optimum: $f_x = 6 - 2 x = 0 \Rightarrow x = 3$;
        $f_y = 4 - 2 y = 0 \Rightarrow y = 2$. Both $\geq 0$ \checkmark,
        so the non-negativity constraints are slack. KKT multipliers
        $\mu_x = \mu_y = 0$. $f^* = 18 + 8 - 9 - 4 = 13$. (A non-binding
        non-negativity constraint in a problem whose unconstrained optimum
        already lies in the positive orthant.)

  \item \emph{Convexity check.} The Hessian of $f$ is
        $\begin{pmatrix} -2 & 0 \\ 0 & -2 \end{pmatrix}$, negative
        definite, so $f$ is strictly concave; the constraint functions
        $x + 2y$, $-x$, $-y$ are linear, hence convex. This is a convex
        programme: any KKT point is automatically the global maximum, and
        no second-order (bordered-Hessian) check is needed.

        \emph{Case A (budget slack, $\lambda = 0$):} stationarity gives the
        unconstrained maximum $(2, 1)$, but $2 + 2 \cdot 1 = 4 > 2$:
        infeasible. Discard.

        \emph{Case B (budget binds, $x, y > 0$):}
        $\mathcal{L} = -(x - 2)^2 - (y - 1)^2 + \lambda(2 - x - 2 y)$.
        FOCs:
        \[
          -2(x - 2) - \lambda = 0
          \;\Rightarrow\; \lambda = -2(x - 2), \qquad
          -2(y - 1) - 2 \lambda = 0
          \;\Rightarrow\; \lambda = -(y - 1) .
        \]
        Equating: $-2(x - 2) = -(y - 1) \Rightarrow 4 - 2 x = 1 - y
        \Rightarrow y = 2 x - 3$. Budget:
        $x + 2(2 x - 3) = 2 \Rightarrow 5 x = 8 \Rightarrow x^* = 8/5$,
        $y^* = 1/5$. Multiplier:
        $\lambda^* = -2(8/5 - 2) = 4/5 \geq 0$ \checkmark (and indeed
        $-(1/5 - 1) = 4/5$, consistent). Non-negativity slack
        ($x^*, y^* > 0$), so $\mu_x = \mu_y = 0$. All KKT conditions hold;
        by sufficiency this is the global maximum:
        \[
          f^* = -\Bigl(\tfrac{8}{5} - 2\Bigr)^2
                - \Bigl(\tfrac{1}{5} - 1\Bigr)^2
              = -\tfrac{4}{25} - \tfrac{16}{25} = -\tfrac{4}{5} .
        \]
        (Geometrically: $(8/5, 1/5)$ is the orthogonal projection of the
        unconstrained peak $(2, 1)$ onto the budget line $x + 2y = 2$.)

  \item \emph{Convexity check.} $4 \ln(1 + x)$ is concave in $x$ (second
        derivative $-4/(1 + x)^2 < 0$) and $y$ enters linearly, so $U$ is
        concave; all constraints are linear. A convex programme: the first
        KKT point found is the global maximum. Since
        $\pd{U}{x} = 4/(1 + x) > 0$ and $\pd{U}{y} = 1 > 0$, utility is
        strictly increasing, so the budget must bind ($\lambda > 0$).
        Lagrangian:
        $\mathcal{L} = 4 \ln(1 + x) + y + \lambda(2 - x - y)
        + \mu_x x + \mu_y y$.

        \emph{Case A ($x, y > 0$, budget binds):} $\mu_x = \mu_y = 0$.
        FOC on $y$: $1 - \lambda = 0 \Rightarrow \lambda = 1$. FOC on $x$:
        $4/(1 + x) = 1 \Rightarrow x = 3$, hence $y = 2 - 3 = -1 < 0$.
        Infeasible; discard.

        \emph{Case B ($y = 0$, $x > 0$, budget binds):} $x^* = 2$. FOC on
        $x$ ($\mu_x = 0$): $\lambda = 4/(1 + 2) = 4/3$. FOC on $y$:
        $1 - \lambda + \mu_y = 0 \Rightarrow \mu_y = 4/3 - 1 = 1/3 \geq 0$
        \checkmark, and $\mu_y\,y = 0$ holds since $y = 0$. All KKT
        conditions satisfied.

        \emph{Case C ($x = 0$, $y > 0$, budget binds):} $y = 2$. FOC on
        $y$: $\lambda = 1$. FOC on $x$:
        $4/(1 + 0) - \lambda + \mu_x = 0 \Rightarrow \mu_x = 1 - 4 = -3
        < 0$: violates dual feasibility. Discard.

        By sufficiency, $(x^*, y^*) = (2, 0)$ is the global maximum with
        $U^* = 4 \ln 3 \approx 4.39$.

  \item (a) The Hessian of $\pi$ is
        $\begin{pmatrix} -2 & 0 \\ 0 & -2 \end{pmatrix}$, negative
        definite: $\pi$ is strictly concave, the constraints are linear,
        so this is a convex programme. Unconstrained optimum:
        $12 - 2 x = 0 \Rightarrow x = 6$,
        $8 - 2 y = 0 \Rightarrow y = 4$; but $6 + 4 = 10 > 8$, infeasible,
        so the capacity constraint binds.
        $\mathcal{L} = 12 x - x^2 + 8 y - y^2 + \lambda(8 - x - y)$. FOCs:
        \[
          12 - 2 x = \lambda, \qquad 8 - 2 y = \lambda
          \;\Rightarrow\; 12 - 2 x = 8 - 2 y
          \;\Rightarrow\; x = y + 2 .
        \]
        Capacity: $(y + 2) + y = 8 \Rightarrow y^* = 3$, $x^* = 5$.
        Multiplier: $\lambda^* = 12 - 10 = 2 \geq 0$ \checkmark (check:
        $8 - 6 = 2$ \checkmark). Both outputs positive, non-negativity
        slack. By convex-programme sufficiency this is the global maximum,
        with no second-order check:
        \[
          \pi^* = 60 - 25 + 24 - 9 = 50 .
        \]
        An extra unit of capacity is worth $\lambda^* = 2$.

        (b) With capacity $12$ the unconstrained optimum $(6, 4)$ is
        feasible: $6 + 4 = 10 \leq 12$. The constraint is slack, so
        complementary slackness $\lambda(12 - 10) = 0$ forces
        $\lambda^* = 0$: extra capacity is worthless. Stationarity reduces
        to the unconstrained FOCs, $(x^*, y^*) = (6, 4)$, and
        \[
          \pi^* = 72 - 36 + 32 - 16 = 52 .
        \]

  \item KKT point: $\mathcal{L}(x, \lambda) = -x^2 + 6 x + \lambda(2 - x)$.
        The unconstrained maximum $x = 3$ violates $x \leq 2$, so the
        constraint binds: $x^* = 2$. Stationarity:
        $-2 x + 6 - \lambda = 0$ at $x = 2$ gives
        $\lambda^* = 2 \geq 0$ \checkmark, and $f^* = -4 + 12 = 8$.

        \emph{Saddle-point verification.} Left inequality: for all
        $x \in \R$,
        \[
          \mathcal{L}(x, 2) = -x^2 + 6 x + 2(2 - x)
                            = -x^2 + 4 x + 4
                            = -(x - 2)^2 + 8 \;\leq\; 8
                            = \mathcal{L}(2, 2),
        \]
        with equality exactly at $x = 2$: so $x^* = 2$ maximises
        $\mathcal{L}(\cdot, \lambda^*)$. Right inequality: for all
        $\lambda \geq 0$,
        \[
          \mathcal{L}(2, \lambda) = -4 + 12 + \lambda(2 - 2) = 8
                                  = \mathcal{L}(2, 2),
        \]
        constant in $\lambda$ (the constraint binds), so
        $\mathcal{L}(2, 2) \leq \mathcal{L}(2, \lambda)$ holds with
        equality. Hence $(x^*, \lambda^*) = (2, 2)$ is a saddle point, and
        by the saddle-point theorem (Problem~20) $x^* = 2$ is the global
        constrained maximum: no concavity or constraint-qualification
        argument was needed.

  \item (i) \emph{Feasibility.} The right inequality reads
        \[
          f(\mathbf{x}^*) + \lambda^*\bigl(b - g(\mathbf{x}^*)\bigr)
          \;\leq\;
          f(\mathbf{x}^*) + \lambda\bigl(b - g(\mathbf{x}^*)\bigr)
          \quad \text{for all } \lambda \geq 0,
        \]
        i.e.\ $(\lambda - \lambda^*)\bigl(b - g(\mathbf{x}^*)\bigr) \geq 0$
        for all $\lambda \geq 0$. If $b - g(\mathbf{x}^*) < 0$, letting
        $\lambda \to \infty$ makes the left side tend to $-\infty$, a
        contradiction. Hence $g(\mathbf{x}^*) \leq b$: $\mathbf{x}^*$ is
        feasible.

        (ii) \emph{Complementary slackness.} Setting $\lambda = 0$ in the
        same inequality gives
        $-\lambda^*\bigl(b - g(\mathbf{x}^*)\bigr) \geq 0$, i.e.\
        $\lambda^*\bigl(b - g(\mathbf{x}^*)\bigr) \leq 0$. But
        $\lambda^* \geq 0$ and, by (i), $b - g(\mathbf{x}^*) \geq 0$, so
        the product is also $\geq 0$. Therefore
        $\lambda^*\bigl(b - g(\mathbf{x}^*)\bigr) = 0$.

        (iii) \emph{Optimality.} Let $\mathbf{x}$ be any feasible point,
        so $b - g(\mathbf{x}) \geq 0$. Using $\lambda^* \geq 0$, the left
        saddle inequality, and (ii) in turn:
        \[
          f(\mathbf{x})
          \;\leq\; f(\mathbf{x})
             + \lambda^*\bigl(b - g(\mathbf{x})\bigr)
          = \mathcal{L}(\mathbf{x}, \lambda^*)
          \;\leq\; \mathcal{L}(\mathbf{x}^*, \lambda^*)
          = f(\mathbf{x}^*)
             + \lambda^*\bigl(b - g(\mathbf{x}^*)\bigr)
          = f(\mathbf{x}^*) .
        \]
        Hence $f(\mathbf{x}^*) \geq f(\mathbf{x})$ for every feasible
        $\mathbf{x}$: $\mathbf{x}^*$ is a global maximum. Nowhere did the
        argument use differentiability, concavity, or a constraint
        qualification; with several constraints the same steps apply to
        each multiplier $\lambda_i$ separately.
\end{solutions}

\section{Additional Problems}

Problems 1--13 repeat the case analysis. Problems 14--58 drill feasibility
and regularity, the graphical method, mixed equality and inequality
constraints, the test of given points, convexity, three variables, economic
programmes with several constraints, standard forms, and non-convex
programmes.

\begin{problems}
  \item Maximise $f(x, y) = x y$ on $x + 2 y \leq 12$, $x \geq 0$,
        $y \geq 0$.
  \item Maximise $f(x, y) = 3 x + 5 y$ on $x + y \leq 8$,
        $2 x + y \leq 10$, $x, y \geq 0$.
  \item Maximise $U(x, y) = \ln(x) + y$ on $x + y \leq 4$, $x, y \geq 0$.
        Is the budget tight?
  \item Minimise $f(x, y) = (x - 1)^2 + (y - 5)^2$ on $x + y \leq 4$,
        $x, y \geq 0$.
  \item Maximise $f(x, y) = x + 3 y$ on $x^2 + y^2 \leq 10$.
  \item Maximise $U = x^{1/2} y^{1/2}$ on $x + y \leq 16$, $x \geq 2$,
        $y \geq 2$. Are the lower bounds binding?
  \item Minimise $f(x) = x^2 + 6 x + 12$ on $x \geq -1$.
  \item Maximise $U = 2 x + \ln(1 + y)$ on $x + y \leq 6$, $x, y \geq 0$.
  \item Maximise $f(x, y) = x y - x^2 - y^2$ on $x + y \leq 4$,
        $x, y \geq 0$.
  \item Maximise $f(x, y) = 4 x + y$ on $x^2 + 4 y^2 \leq 8$,
        $y \leq 1$.

  \item Using the figure of the Core Rules: for each of
        $P_{1}=(3,1)$, $P_{2}=(0,1)$, $P_{3}=(1,4)$ and $P_{4}=(4,1)$,
        decide whether the point lies in the feasible set $C$, and for
        those that do, state which constraints are binding and which are
        slack.
  \item Still from the figure: at a point where only the constraint
        $x\le3$ binds, which Kuhn--Tucker multipliers must be zero, and
        what does the remaining stationarity condition say?
  \item Sketch the gradient of the objective $f(x,y)=x+y$ on the same
        figure and read off, without computing, at which corner of $C$ it
        is maximised. Then confirm the answer with the Kuhn--Tucker
        conditions.

  \item Let $S=\{(x,y): y\ge(x-1)^{2},\ y\le3,\ x+y\le4\}$. For each of
        $(1,0)$, $(0,1)$, $(2,2)$, $(1,3)$, $(3,1)$, $(-1,3)$ and
        $(1-\sqrt3,\,3)$ decide whether it lies in $S$, and for the feasible
        ones name the binding and the slack constraints.
  \item Let $S=\{(x,y): x^{2}+y^{2}\le25,\ x-y\le1,\ y\le4\}$. Classify the
        points $(3,4)$, $(4,3)$, $(0,0)$, $(5,0)$, $(-3,4)$, $(1,5)$,
        $(-5,0)$ in the same way.
  \item (Regularity.) Show that the programme with feasible set
        $\{x^{2}+y^{2}\le4,\ y\ge0,\ x\le1\}$ is regular, (a) by Slater,
        (b) by checking LICQ at every corner and on every boundary piece.
  \item Is the programme with feasible set $\{y\le4-x^{2},\ y\ge0\}$
        regular? Check the corners where the parabola meets the axis.
  \item (Cusp.) Let $S=\{(x,y): y\le(2-x)^{3},\ y\ge0\}$. (a) Sketch $S$
        near $(2,0)$ and show that $S$ lies in $x\le2$. (b) Show that LICQ
        fails at $(2,0)$. (c) Show that $(2,0)$ maximises $f=x$ on $S$ but
        is not a KKT point.
  \item Show that the programme in $\R^{3}$ with constraints
        $x^{2}+y^{2}+z^{2}\le9$, $x+y+z\ge3$, $z\ge0$ is regular.
  \item The set $\{x^{2}+y^{2}\ge1,\ 0\le x\le2,\ 0\le y\le2\}$ (a square
        with a quarter disc removed) is not convex. Show that the programme
        is nevertheless regular.
  \item (Graphical.) Let $S=\{(x,y): x^{2}\le y\le2\}$ and $f=x+y$. Sketch
        $S$ and the level lines, find the minimum and the maximum of $f$ on
        $S$ graphically, and confirm each with the KKT conditions, including
        the multiplier signs.
  \item On the upper half annulus $\{1\le x^{2}+y^{2}\le4,\ y\ge0\}$ find
        the minimum and the maximum of $f=x+2y$ graphically and confirm
        with KKT. Is the feasible set convex?
  \item On $\{xy\ge4,\ x+y\le5,\ x>0\}$ find the minimum and the maximum
        of $f=x^{2}+y^{2}$ graphically (level circles) and confirm both with
        KKT. How many maximisers are there?
  \item On the disc $x^{2}+y^{2}\le4$ find the maximum and the minimum of
        $f=y-x^{2}$ (level curves are parabolas). Find \emph{all} KKT points
        of the minimisation and explain why one of them is not a minimiser.
  \item Let $S=\{(x,y): 0\le x\le2,\ 0\le y\le x^{2}\}$, the region under a
        convex parabola, and minimise the squared distance
        $f=x^{2}+(y-3)^{2}$ to $(0,3)$. (a) Is this a convex programme?
        (b) Find all KKT points. (c) Show that at $(0,0)$ the multipliers
        are not unique, and explain why. (d) Which KKT point is the
        minimiser?
  \item Minimise $f=xy$ on $\{x+y\le4,\ x\ge0,\ y\ge0\}$. Show that the set
        of minimisers is a union of two segments, and that every point of it
        satisfies the KKT conditions.
  \item (Testing points.) For $\max f=x+2y$ s.t.\ $x^{2}+y^{2}\le5$,
        $y\le1$, test each of $P=(1,2)$, $Q=(2,1)$, $R=(0,1)$, $S=(-2,1)$,
        $T=(0,0)$: feasible? active set? KKT point (with multipliers)? Is the
        surviving point optimal?
  \item For $\max f=\ln(1+x)+\ln(1+y)+z$ s.t.\ $x+y+z\le4$,
        $x,y,z\ge0$, test the points $(1,1,2)$, $(0,0,4)$, $(2,2,0)$ and
        $(1,1,1)$. Does optimality follow for the KKT point?
  \item For $\max f=x^{2}+y^{2}$ s.t.\ $x+2y\le4$, $x,y\ge0$, show that
        $(4,0)$, $(0,2)$, $(0,0)$ and $(0.8,1.6)$ all satisfy the KKT
        conditions. Which is optimal, and why does KKT not decide it here?
  \item (Equality and inequality.) Maximise $f=xy-\tfrac12y^{2}$ s.t.\
        $x+y=6$, $y\ge3$, $x,y\ge0$. (a) Solve with a free multiplier for
        the equality. (b) Redo with the equality as two inequalities and
        show that the two multipliers are not unique. (c) Interpret the
        multiplier of $y\ge3$.
  \item Minimise $x^{2}+y^{2}+z^{2}$ s.t.\ $x+y+z=3$ and $x\ge2$.
  \item Maximise $\ln x+\ln y$ s.t.\ $x+2y=10$ and $x\ge6$.
  \item Maximise $f=-(x-1)^{2}-(y-1)^{2}$ s.t.\ $x+y=4$, $y\le1.5$,
        $x,y\ge0$. Why is a negative multiplier for the equality not an
        error?
  \item (Min with $\ge$, then max.) Write the full KKT system, in folded
        form, for $\min f=x^{2}+2y^{2}$ s.t.\ $x+y\ge3$, $x,y\ge0$, and solve
        it. Then write the system for the maximisation of the same $f$ on
        the same set and show that it has no solution. What does that prove?
  \item For $f=(x-1)^{2}+(y-2)^{2}$ on $\{x+2y\ge6,\ x\le3,\ y\le3\}$:
        write the KKT system for the minimum and solve it; then write the
        system for the maximum and verify it at the optimal corner.
  \item (Convex or not?) Decide which programmes are convex:
        (a) $\min(x-2)^{2}+(y-1)^{2}$ s.t.\ $x^{2}+y^{2}\le1$;
        (b) $\min x+y$ s.t.\ $x^{2}+y^{2}\ge1$;
        (c) $\min x^{2}+y^{2}$ s.t.\ $x^{2}-y^{2}\ge1$;
        (d) $\min -x-y$ s.t.\ $y\le x^{2}$, $0\le x\le2$, $y\ge0$;
        (e) $\min x^{4}+x^{2}-2xy+y^{2}$ s.t.\ $x+y\ge3$;
        (f) $\max x^{2}+y^{2}$ s.t.\ $x^{2}+y^{2}\le1$.
  \item Solve programme (c) of the previous problem. How many minimisers
        are there?
  \item Solve programme (e) of Additional Problem~36, after checking with
        the Hessian that the objective is convex.
  \item (Three variables.) Minimise $(x-2)^{2}+(y-2)^{2}+(z-2)^{2}$ s.t.\
        $x+y+z\le3$ and $x+2y\le2$ by case analysis, and certify the
        answer.
  \item Maximise $f=xy+yz+zx$ on the ball $x^{2}+y^{2}+z^{2}\le3$ (no sign
        restrictions). Use symmetry to reduce the cases. Is $f$ concave?
        How many maximisers are there?
  \item Maximise $U=x+2\sqrt{y}+\ln(1+z)$ s.t.\ $x+y+z\le1.5$, $x\ge1$,
        $y,z\ge0$.
  \item (Two budgets.) A consumer with $U=x^{1/2}y^{1/2}$ has $40$ euros,
        prices $p_x=2$, $p_y=1$, and $40$ ration coupons; each unit of $x$
        needs one coupon, each unit of $y$ two. (a) Formulate with
        inequalities. (b) Show that the optimum of either budget alone
        violates the other. (c) Verify the point where both bind and
        interpret the multipliers. Would the consumer pay $1.50$ euros for an
        extra coupon?
  \item A firm makes products $x,y$ with profit
        $\pi=20x-x^{2}+16y-y^{2}$, using raw material $A$ ($x+y\le10$) and
        $B$ ($2x+y\le15$). Show that the unconstrained optimum and the
        single-constraint optima are infeasible, find the optimum where both
        bind, and interpret $\lambda_A$, $\lambda_B$.
  \item (Budget left unspent.) Advertising $a_1,a_2\ge0$ in two channels
        yields profit
        $\pi=60(1-\mathrm{e}^{-a_1/20})+40(1-\mathrm{e}^{-a_2/10})-a_1-a_2$,
        and at most $B$ may be spent. Solve for $B=50$ and for $B=30$, and
        interpret the multiplier in each case.
  \item Maximise and minimise $f=xy$ s.t.\ $x+y=10$ and $2\le x\le6$, with
        the KKT conditions. Which bound binds at which optimum?
  \item (Standard form.) Rewrite
        $\max xy-z^{2}+2z$ s.t.\ $x^{2}+y\le5$, $x-y\ge-1$, $y\ge0$,
        $z\le0$, $x$ free, (a) as ``$\max F$ s.t.\ $G_i\ge0$, all variables
        $\ge0$'', (b) as ``$\min F$ s.t.\ $G_i\le0$, all variables $\ge0$''.
  \item Rewrite $\min(x-1)^{2}+\mathrm{e}^{w}$ s.t.\ $x+w=2$, $w\le0$, $x$
        free, in both target forms, and solve the problem.
  \item (Lagrange versus KKT.) A wire of length $20$ is cut; a piece of
        length $x$ is bent into a square, the rest $y=20-x$ into a circle.
        (a) Minimise and maximise the total area by substitution on
        $[0,20]$. (b) Solve with Lagrange on $x+y=20$ and classify with the
        bordered Hessian. (c) Use KKT on $\{x+y=20,\ x,y\ge0\}$ to explain
        why Lagrange finds the minimum but misses the maximum.
  \item Show that $\max x^{2}+xy+y^{2}$ s.t.\ $x+y\ge2$, $x,y\ge0$ has no
        solution, by showing that no point satisfies the KKT conditions.
        Then solve the minimisation.
  \item (Modelling.) A firm will place a warehouse at a point $p$ inside a
        zoning area, the disc $\|p\|\le5$, and wants to minimise the sum of
        squared distances to two shops at $(6,0)$ and $(6,4)$.
        (a) Name the objective and say what the constraint function
        describes. (b) Is the feasible set convex? Is the programme convex,
        regular? (c) Without computing: are the KKT conditions necessary,
        sufficient? If Lagrange gives a point on the boundary circle with a
        non-negative multiplier, what follows? (d) Find the optimum.
  \item For each programme decide whether the KKT conditions are necessary
        and whether they are sufficient, with a reason:
        (a) $\max\ln x+\ln y$ s.t.\ $x+y\le10$;
        (b) $\max x^{2}+y^{2}$ s.t.\ $x+y\le1$, $x,y\ge0$;
        (c) $\max x$ s.t.\ $y\le(1-x)^{3}$, $y\ge0$;
        (d) $\min x^{2}+y^{2}$ s.t.\ $x^{2}+y^{2}\ge1$.
  \item Explain in a few lines why the multipliers at a regular optimum are
        unique, and what non-unique multipliers signal. Give an example of
        each from this chapter.
  \item (Portfolio, Kuhn--Tucker.) Three uncorrelated assets have expected
        returns $8\%$, $5\%$, $3\%$ and variances $100$, $36$, $16$. With
        shares $x_1,x_2$ and $x_3=1-x_1-x_2$, minimise the variance subject
        to an expected return of at least $6\%$, without short sales
        ($x_1,x_2\ge0$, $x_1+x_2\le1$). (a) Write the programme and its KKT
        conditions. (b) Taking the return constraint as binding and the
        others as slack, solve and verify every condition. (c) Argue global
        optimality.
  \item Repeat the previous problem with a target return of $7.5\%$. Show
        that the guess ``only the return constraint binds'' now fails, find
        the correct active set and verify all conditions.
  \item (Non-convex profit.) A firm earns $\pi=6\ln(1+x_1x_2)-x_1-x_2$
        from inputs $x_1,x_2\ge0$ and may spend at most $B$, i.e.\
        $x_1+x_2\le B$. (a) Find the interior critical points of $\pi$ and
        classify them. (b) For $B=20$ show that the saddle point satisfies
        all KKT conditions with $\lambda=0$, yet is not optimal; find the
        optimum. (c) Solve for $B=4$ and for $B=0.5$, interpreting
        $\lambda$ in each case.
  \item (Revenue instead of profit.) Maximise revenue $R=6\ln(1+x_1x_2)$
        s.t.\ $x_1+x_2\le B$, $x_1,x_2\ge0$. Why does the budget always
        bind, so that it may be written as an equality? Give the optimum and
        $\lambda$ for $B=4$ and $B=20$ and compare with the previous
        problem.
  \item (Domain as constraint.) Maximise $f=x+2y+\sqrt{5-x^{2}-y^{2}}$
        over its natural domain, with no sign restrictions. Write the domain
        as an explicit inequality constraint and show that the maximiser
        satisfies the KKT conditions with a slack domain constraint.
  \item Consider $\max f=x+y+(y-x^{2})^{3/2}$ on its domain $y\ge x^{2}$,
        $x$ and $y$ unrestricted. Show that no point satisfies the KKT
        conditions and conclude that no maximum exists.
\end{problems}

\section{Additional Solutions}

\begin{solutions}
  \item Interior solution with the budget binding.
        $\mathcal{L} = x y + \lambda(12 - x - 2 y)$. FOCs: $y = \lambda$;
        $x = 2 \lambda \Rightarrow x = 2 y$. Constraint:
        $2 y + 2 y = 4 y = 12 \Rightarrow y^* = 3$, $x^* = 6$.
        $\lambda^* = 3 \geq 0$ \checkmark. Non-negativity slack.
        $f^* = 18$.

  \item Linear programming. Vertices of the feasible polygon (intersect
        $x + y = 8$, $2 x + y = 10$, $x = 0$, $y = 0$): $(0, 0)$,
        $(5, 0)$, $(2, 6)$, $(0, 8)$. (Check: $(5, 0)$: $5 \leq 8$,
        $10 \leq 10$. $(2, 6)$: $8 \leq 8$, $10 \leq 10$. $(0, 8)$:
        $8 \leq 8$, $8 \leq 10$.) Evaluate $f = 3 x + 5 y$: $0$, $15$,
        $36$, $40$. \textbf{Maximum} at $(0, 8)$ with $f^* = 40$. Only
        $x + y \leq 8$ binds, plus $x \geq 0$. Multipliers: $\lambda_1$ on
        the first budget, $\mu_x$ on $x \geq 0$. FOCs:
        $3 - \lambda_1 + \mu_x = 0 \Rightarrow \mu_x = \lambda_1 - 3$;
        $5 - \lambda_1 = 0 \Rightarrow \lambda_1 = 5$. Then
        $\mu_x = 2 \geq 0$ \checkmark. The second budget is slack with
        multiplier $0$.

  \item Try an interior solution with the budget binding.
        $\mathcal{L} = \ln x + y + \lambda(4 - x - y)$. FOCs:
        $1/x = \lambda$; $1 = \lambda \Rightarrow x^* = 1$,
        $\lambda^* = 1$. Then $y^* = 3 \geq 0$ \checkmark. All
        non-negativity constraints slack with multipliers $0$.
        $U^* = \ln 1 + 3 = 3$. The budget binds (with
        $\lambda = 1 > 0$).

  \item The unconstrained minimum at $(1, 5)$ has $1 + 5 = 6 > 4$:
        infeasible. Try the budget binding in the interior:
        $\mathcal{L} = (x - 1)^2 + (y - 5)^2 + \mu(x + y - 4)$ with
        $\mu \geq 0$ (min form with $x + y - 4 \leq 0$). FOCs:
        \[
          2(x - 1) + \mu = 0, \qquad 2(y - 5) + \mu = 0
          \;\Rightarrow\; x - 1 = y - 5 \;\Rightarrow\; y = x + 4 .
        \]
        Constraint: $x + x + 4 = 4 \Rightarrow x^* = 0$, $y^* = 4$. But
        $x^* = 0$ lies on the boundary of $x \geq 0$. Check:
        $\mu = -2(x - 1) = 2 \geq 0$ \checkmark. Non-negativity
        $x \geq 0$ also binds at this corner; let $\mu_x$ be its
        multiplier. Modified FOC:
        $2(x - 1) + \mu - \mu_x = 0 \Rightarrow -2 + 2 - \mu_x = 0
        \Rightarrow \mu_x = 0$. So the lower-bound constraint binds with
        zero multiplier (a boundary solution with non-strict KKT).
        $f^* = 1 + 1 = 2$.

  \item The circle constraint should bind.
        $\mathcal{L} = x + 3 y + \lambda(10 - x^2 - y^2)$. FOCs:
        $1 = 2 \lambda x$; $3 = 2 \lambda y \Rightarrow
        x = 1/(2 \lambda)$, $y = 3/(2 \lambda) = 3 x$. Constraint:
        $x^2 + 9 x^2 = 10 x^2 = 10 \Rightarrow x^* = 1$ (positive for a
        max), $y^* = 3$. $\lambda^* = 1/(2 x^*) = 1/2 \geq 0$ \checkmark.
        $f^* = 1 + 9 = 10$.

  \item Try an interior solution with the budget binding and the lower
        bounds slack: the FOC for $U = \sqrt{x y}$ gives $y/x = 1$
        (symmetric Cobb--Douglas), so $x = y$. Budget:
        $2 x = 16 \Rightarrow x = y = 8 > 2$ \checkmark. The lower bounds
        are slack with multipliers $0$. For $\lambda^*$ compute:
        $\pd{U}{x} = \tfrac12 \sqrt{y/x} = \tfrac12$ at $x = y = 8$, and
        the FOC sets this equal to $\lambda \cdot 1$, so
        $\lambda^* = 1/2$. $U^* = 8$.

  \item Constraint $x \geq -1$, i.e.\ $-1 - x \leq 0$. Unconstrained
        minimum: $f'(x) = 2 x + 6 = 0 \Rightarrow x = -3$, infeasible
        (violates $x \geq -1$). So the lower bound binds: $x^* = -1$.
        $f(-1) = 1 - 6 + 12 = 7$. KKT for a min: $f'(x) - \mu = 0$ where
        $\mu \geq 0$ is the multiplier on $-1 - x \leq 0$.
        $f'(-1) = 4 = \mu \geq 0$ \checkmark.

  \item Try an interior solution with the budget binding.
        $\mathcal{L} = 2 x + \ln(1 + y) + \lambda(6 - x - y)$. FOCs:
        $2 = \lambda$; $1/(1 + y) = \lambda = 2 \Rightarrow
        1 + y = 1/2 \Rightarrow y = -1/2$, infeasible. So the
        $y$-non-negativity binds: $y^* = 0$. Then $x^* = 6$. Check the FOC
        at $y = 0$: $1/(1 + 0) = 1$ versus $\lambda = 2$. Multiplier on
        $y \geq 0$: $\mu_y = \lambda - 1 = 1 \geq 0$ \checkmark.
        $\lambda^* = 2$. $U^* = 12 + 0 = 12$.

  \item Unconstrained maximum of $f = x y - x^2 - y^2$:
        $f_x = y - 2 x = 0$, $f_y = x - 2 y = 0 \Rightarrow x = y = 0$
        (and $f^* = 0$). At $(0, 0)$ all constraints are satisfied (budget
        slack with multiplier $0$, non-negativity binding with multipliers
        $0$). Hessian
        $\begin{pmatrix} -2 & 1 \\ 1 & -2 \end{pmatrix}$ with leading
        principal minors $D_1 = -2 < 0$, $D_2 = 3 > 0$: negative definite
        $\Rightarrow$ maximum. $f^* = 0$. (A trivial KKT problem: the
        unconstrained optimum is at the origin, which is feasible.)

  \item \emph{Case A (ellipse and $y \leq 1$ both binding):} $y = 1$,
        $x^2 + 4 = 8 \Rightarrow x^2 = 4 \Rightarrow x = 2$ (positive for
        a max), $f = 8 + 1 = 9$.
        $\mathcal{L} = 4 x + y + \lambda_1(8 - x^2 - 4 y^2)
        + \lambda_2(1 - y)$. FOCs:
        $4 = 2 \lambda_1 x \Rightarrow \lambda_1 = 1$;
        $1 = 8 \lambda_1 y + \lambda_2 = 8 + \lambda_2 \Rightarrow
        \lambda_2 = -7 < 0$: this combination fails.

        \emph{Case B (only the ellipse binding, $y < 1$):}
        $4 = 2 \lambda_1 x$, $1 = 8 \lambda_1 y \Rightarrow
        x = 2/\lambda_1$, $y = 1/(8 \lambda_1)$. Constraint:
        \[
          \frac{4}{\lambda_1^2} + \frac{4}{64 \lambda_1^2}
          = \frac{4}{\lambda_1^2} + \frac{1}{16 \lambda_1^2}
          = \frac{64 + 1}{16 \lambda_1^2}
          = \frac{65}{16 \lambda_1^2} = 8
          \;\Rightarrow\; \lambda_1^2 = \frac{65}{128}
          \;\Rightarrow\; \lambda_1 = \sqrt{65/128} .
        \]
        Then $y = 1/(8 \lambda_1) = \sqrt{128/65}/8 = \sqrt{2/65}$. Check
        $y \leq 1$: $\sqrt{2/65} \approx 0.175 \leq 1$ \checkmark. So
        $x^* = 2 \sqrt{128/65}$, $y^* = \sqrt{2/65}$, and
        \[
          f^* = 4 x^* + y^*
              = 8 \sqrt{128/65} + \sqrt{2/65}
              = \frac{8 \cdot 8 \sqrt{2} + \sqrt{2}}{\sqrt{65}}
              = \frac{65 \sqrt{2}}{\sqrt{65}}
              = \sqrt{65 \cdot 2} = \sqrt{130} \approx 11.4 .
        \]
        $\lambda_2 = 0$ in this case.

  \item $P_{1}=(3,1)$: $(3-1)^{2}+0=4\le9$ (slack), $x=3$
        \textbf{binding}, $x+y=4\ge1$ (slack). Feasible, one binding
        constraint. \;$P_{2}=(0,1)$: $1\le9$ slack, $x=0<3$ slack,
        $x+y=1$ \textbf{binding}. Feasible. \;$P_{3}=(1,4)$:
        $0+9=9$ so the circle \textbf{binds}, $x=1<3$ slack, $x+y=5$
        slack. Feasible. \;$P_{4}=(4,1)$: $x=4>3$ violates the second
        constraint --- \textbf{not feasible} (the circle constraint alone
        would have allowed it, since $9+0=9$).
  \item Complementary slackness forces the multipliers of the two slack
        constraints to vanish: the circle's multiplier and the
        multiplier on $x+y\ge1$ are both $0$. Stationarity then reduces to
        $\grad f=\mu\,\grad(x)= \mu\,(1,0)^{T}$, i.e. $f_{y}=0$ and
        $f_{x}=\mu\ge0$: the objective must be flat in $y$ and
        non-decreasing in $x$ at that point, which is what makes pushing
        against the wall $x=3$ optimal.
  \item $\grad f=(1,1)^{T}$, a constant arrow pointing up-right, so the
        level lines $x+y=\text{const}$ are the anti-diagonals and the
        optimum is the feasible point furthest up-right. That is the corner
        where $x=3$ meets the circle: $x=3$ and
        $(3-1)^{2}+(y-1)^{2}=9$ give $(y-1)^{2}=5$, so
        $y=1+\sqrt5\approx3.236$ and $f=4+\sqrt5\approx6.236$. Both
        constraints bind there, so two multipliers may be positive while
        the multiplier on $x+y\ge1$ is zero; solving the stationarity
        conditions confirms non-negative multipliers, so the Kuhn--Tucker
        conditions hold.

  \item Constraints: (i) $y\ge(x-1)^{2}$, (ii) $y\le3$, (iii) $x+y\le4$.
        $(1,0)$: $0\ge0$ \textbf{binds}, (ii), (iii) slack. $(0,1)$:
        $1\ge1$ \textbf{binds}, others slack. $(2,2)$: (i) $2\ge1$ slack,
        (ii) slack, $2+2=4$ \textbf{binds}. $(1,3)$: (i) slack, (ii) and
        (iii) \textbf{both bind}, a corner. $(3,1)$: $1\ge4$ fails,
        \textbf{infeasible}. $(-1,3)$: $3\ge4$ fails, \textbf{infeasible}.
        $(1-\sqrt3,3)$: $(-\sqrt3)^{2}=3\le3$, so (i) and (ii) \textbf{both
        bind}; $x+y=4-\sqrt3\approx2.27$, (iii) slack: the left corner where
        the parabola meets $y=3$.
  \item Constraints: (i) $x^{2}+y^{2}\le25$, (ii) $x-y\le1$, (iii)
        $y\le4$. $(3,4)$: (i) and (iii) bind, (ii) slack ($-1$). $(4,3)$:
        (i) and (ii) bind, (iii) slack. $(0,0)$: all slack, an interior
        point. $(5,0)$: $x-y=5>1$, \textbf{infeasible}. $(-3,4)$: (i) and
        (iii) bind, (ii) slack. $(1,5)$: $26>25$ (and $5>4$),
        \textbf{infeasible}. $(-5,0)$: (i) binds, (ii), (iii) slack.
  \item (a) $x^{2}+y^{2}$ is convex, $-y$ and $x$ are linear, so the
        constraint set is convex, and $(0,1)$ satisfies every inequality
        strictly ($1<4$, $1>0$, $0<1$): Slater holds, the programme is
        regular. (b) Gradients: circle $(2x,2y)$, never zero on the circle;
        $y\ge0$: $(0,-1)$; $x\le1$: $(1,0)$. Corners in the set:
        circle $\cap\ \{y=0\}$ at $(-2,0)$ (the point $(2,0)$ violates
        $x\le1$): $(-4,0)$, $(0,-1)$ independent. Circle $\cap\ \{x=1\}$ at
        $(1,\sqrt3)$: $(2,2\sqrt3)$, $(1,0)$, determinant $-2\sqrt3\neq0$.
        $\{y=0\}\cap\{x=1\}$ at $(1,0)$: $(0,-1)$, $(1,0)$ independent. No
        three constraints meet. LICQ holds everywhere: regular.
  \item $g_1=y+x^{2}-4\le0$ has $\grad g_1=(2x,1)\neq0$; $g_2=-y\le0$ has
        $(0,-1)$. They meet at $(\pm2,0)$, where $(\pm4,1)$ and $(0,-1)$ have
        determinant $\mp4\neq0$: independent. \textbf{Regular} (also by
        Slater: the set is convex and $(0,1)$ is strictly inside).
  \item (a) $y\ge0$ and $y\le(2-x)^{3}$ require $(2-x)^{3}\ge0$, i.e.\
        $x\le2$. Near $(2,0)$ the set is squeezed between $y=0$ and the curve
        $y=(2-x)^{3}$, which touches the axis at $x=2$: a cusp pointing
        right. (b) At $(2,0)$ both constraints are active,
        $\grad\bigl(y-(2-x)^{3}\bigr)=(3(2-x)^{2},1)=(0,1)$ and
        $\grad(-y)=(0,-1)$: parallel, LICQ fails. (c) $x\le2$ on $S$ and
        $(2,0)\in S$, so $(2,0)$ maximises $f=x$. KKT would need
        $(1,0)=\lambda_1(0,1)+\lambda_2(0,-1)$, whose first component reads
        $1=0$: \textbf{not a KKT point}. The optimum escapes the conditions
        because the programme is not regular there.
  \item The ball constraint $x^{2}+y^{2}+z^{2}\le9$ is convex, the others
        are linear, so the feasible set is convex. Slater point
        $(1,1,1.5)$: $1+1+2.25=4.25<9$, $3.5>3$, $1.5>0$. Hence
        \textbf{regular}.
  \item Write $1-x^{2}-y^{2}\le0$, gradient $(-2x,-2y)$, zero only at the
        origin, which is not feasible. The circle meets $x=0$ at $(0,1)$:
        gradients $(0,-2)$ and $(-1,0)$, independent; it meets $y=0$ at
        $(1,0)$: $(-2,0)$ and $(0,-1)$, independent; it does not reach $x=2$
        or $y=2$. The remaining corners $(2,0)$, $(2,2)$, $(0,2)$ involve two
        independent linear constraints. LICQ holds at every feasible point:
        \textbf{regular}, although the set is not convex. Regularity and
        convexity are different questions.
  \item $S$ is the region between the parabola and the line $y=2$, convex.
        Level lines $x+y=k$ move up-right as $k$ grows.
        \emph{Minimum:} the lowest line touching $S$ is tangent to the
        parabola, slope $-1=2x$: $(-\tfrac12,\tfrac14)$, $f=-\tfrac14$. KKT
        (min, $g=x^{2}-y\le0$): $\mathcal{L}=x+y+\lambda(x^{2}-y)$,
        $1+2\lambda x=0$, $1-\lambda=0$: $\lambda=1\ge0$, $x=-\tfrac12$.
        \checkmark\ \emph{Maximum:} the last contact is the right corner
        $(\sqrt2,2)$, $f=2+\sqrt2\approx3.41$. KKT (max, both bind):
        $(1,1)=\lambda_1(0,1)+\lambda_2(2\sqrt2,-1)$ gives
        $\lambda_2=\tfrac{\sqrt2}{4}$, $\lambda_1=1+\tfrac{\sqrt2}4$, both
        positive. \checkmark\ (A linear objective on a convex set: both
        problems are convex programmes, so both KKT points are global.)
  \item The set is not convex: $(-1.5,0)$ and $(1.5,0)$ belong to it, their
        midpoint $(0,0)$ does not. \emph{Maximum:} level lines $x+2y=k$
        last touch the outer arc where it is tangent, at
        $2\cdot\tfrac{(1,2)}{\sqrt5}=(\tfrac2{\sqrt5},\tfrac4{\sqrt5})$,
        $f=2\sqrt5\approx4.47$. KKT: $(1,2)=\lambda(2x,2y)=\lambda(\tfrac4{\sqrt5},\tfrac8{\sqrt5})$,
        $\lambda=\tfrac{\sqrt5}{4}>0$. \emph{Minimum:} the tangency with the
        inner arc would be at $-\tfrac{(1,2)}{\sqrt5}$, which has $y<0$:
        infeasible. The lowest contact is the corner $(-2,0)$, $f=-2$
        (along the arcs $2\cos\theta+4\sin\theta\ge-2$ and
        $\cos\theta+2\sin\theta\ge-1$ for $\theta\in[0,\pi]$). KKT (min,
        $x^{2}+y^{2}\le4$ with $\mu_1$, $-y\le0$ with $\mu_2$):
        $(1,2)+\mu_1(-4,0)+\mu_2(0,-1)=0$, so $\mu_1=\tfrac14$, $\mu_2=2$,
        both $\ge0$. \checkmark
  \item Level circles around the origin. \emph{Minimum:} the smallest
        circle reaching $xy\ge4$ touches the hyperbola at $(2,2)$, $f=8$,
        and $x+y=4\le5$ there. KKT (min, $g=4-xy\le0$):
        $(4,4)+\lambda(-y,-x)=(4,4)-2\lambda(1,1)=0$, $\lambda=2$.
        \emph{Maximum:} the farthest feasible points are the corners where
        $xy=4$ meets $x+y=5$: $(1,4)$ and $(4,1)$, $f=17$, \textbf{two
        maximisers}. KKT at $(4,1)$ (max, both bind):
        $(8,2)=\lambda_1(1,1)+\lambda_2(-1,-4)$, so $8=\lambda_1-\lambda_2$,
        $2=\lambda_1-4\lambda_2$: $\lambda_2=2$, $\lambda_1=10$; at $(1,4)$
        the same by symmetry. (The set is convex, so the minimisation is a
        convex programme; the maximisation of a convex $f$ is not.)
  \item \emph{Maximum:} level curves $y=x^{2}+k$ are upward parabolas;
        raising $k$, the last contact is $(0,2)$, $f=2$. KKT: $(0,1)=\lambda(0,4)$,
        $\lambda=\tfrac14$. $f$ is concave and the disc convex: a convex
        programme, so this is the global maximum. \emph{Minimum:}
        $\mathcal{L}=y-x^{2}+\lambda(x^{2}+y^{2}-4)$, stationarity
        $2x(\lambda-1)=0$ and $1+2\lambda y=0$. $\lambda=0$ is impossible
        ($1=0$), so the circle binds. Either $\lambda=1$, $y=-\tfrac12$,
        $x=\pm\tfrac{\sqrt{15}}{2}$, $f=-\tfrac{17}{4}$; or $x=0$, $y=\pm2$:
        $y=2$ needs $\lambda=-\tfrac14<0$ (out), $y=-2$ gives
        $\lambda=\tfrac14\ge0$, $f=-2$. So the KKT points are
        $(\pm\tfrac{\sqrt{15}}2,-\tfrac12)$ and $(0,-2)$. The minimum is
        $-\tfrac{17}{4}$ at the two symmetric points; $(0,-2)$ is \textbf{not}
        a minimiser. Minimising a concave function is not a convex
        programme: at $(0,-2)$ the first-order conditions hold, but moving
        along the circle ($f=y^{2}+y-4$ there) lowers $f$, a second-order
        effect KKT cannot see.
  \item (a) The objective is convex, but the feasible set is not:
        $(0,0)$ and $(2,4)$ lie in $S$, the midpoint $(1,2)$ violates
        $y\le x^{2}$. \textbf{Not} a convex programme.
        (b) Min form, $\mathcal{L}=f+\lambda_1(y-x^{2})+\mu_1(x-2)-\mu_2x-\mu_3y$.
        Interior: $\grad f=0$ at $(0,3)\notin S$. Only the parabola active:
        $(2x-2\lambda_1x,\ 2(y-3)+\lambda_1)=0$ gives $\lambda_1=1$,
        $y=\tfrac52$, $x=\sqrt{5/2}\approx1.58\in(0,2)$: a KKT point, $f=\tfrac52+\tfrac14=\tfrac{11}4$.
        Only $x\le2$ active (points $(2,y)$, $0<y<4$): $(4,\,2(y-3))+\mu_1(1,0)=0$
        needs $\mu_1=-4$: out. Only $y\ge0$ active with $x>0$:
        $2x=0$, impossible. Corner $(2,4)$: $(4,2)+\lambda_1(-4,1)+\mu_1(1,0)=0$
        gives $\lambda_1=-2$: out. Corner $(2,0)$: $\mu_1=-4$: out. Corner
        $(0,0)$: see (c), a KKT point with $f=9$.
        (c) At $(0,0)$ three constraints are active ($y\le x^{2}$, $x\ge0$,
        $y\ge0$) with gradients $(0,1)$, $(-1,0)$, $(0,-1)$: three vectors
        in the plane, necessarily dependent. Stationarity
        $(0,-6)+\lambda_1(0,1)+\mu_2(-1,0)+\mu_3(0,-1)=0$ gives $\mu_2=0$
        and $\lambda_1-\mu_3=6$: the one-parameter family
        $(\lambda_1,\mu_3)=(6+t,\,t)$, $t\ge0$. Not unique, because LICQ
        fails.
        (d) Compare: $f(0,0)=9>\tfrac{11}4$. The minimiser is
        $(\sqrt{5/2},\tfrac52)$; the corner is a KKT point that is not
        optimal, possible because the programme is not convex.
  \item On the triangle $f=xy\ge0$, and $f=0$ exactly when $x=0$ or
        $y=0$: the minimisers are the segments $\{(0,y):0\le y\le4\}$ and
        $\{(x,0):0\le x\le4\}$. Folded KKT for the minimum with
        $\mathcal{L}=xy+\lambda(x+y-4)$: $\mathcal{L}_x=y+\lambda\ge0$,
        $x\mathcal{L}_x=0$; $\mathcal{L}_y=x+\lambda\ge0$, $y\mathcal{L}_y=0$;
        $\lambda\ge0$, $\lambda(x+y-4)=0$. At $(0,y_0)$ take $\lambda=0$:
        $\mathcal{L}_x=y_0\ge0$ with $x=0$, and $\mathcal{L}_y=0$:
        satisfied (also at $y_0=4$, where the budget binds with
        $\lambda=0$). By symmetry the same on the other axis. A non-convex
        objective with a continuum of optima, all of them KKT points.
  \item Constraints (i) $x^{2}+y^{2}\le5$ ($\lambda_1$), (ii) $y\le1$
        ($\lambda_2$); stationarity $(1,2)=\lambda_1(2x,2y)+\lambda_2(0,1)$.
        $P=(1,2)$: $y=2>1$, \textbf{infeasible}. $Q=(2,1)$: both bind;
        $1=4\lambda_1$, $2=2\lambda_1+\lambda_2$: $\lambda_1=\tfrac14$,
        $\lambda_2=\tfrac32$, both positive: \textbf{KKT point}.
        $R=(0,1)$: only (ii) active, $(1,2)=\lambda_2(0,1)$ impossible:
        not KKT. $S=(-2,1)$: both bind, $1=-4\lambda_1$ gives
        $\lambda_1=-\tfrac14<0$: \textbf{fails the sign test}. $T=(0,0)$:
        nothing active, $\grad f\neq0$: not KKT. The programme is convex
        (linear objective, convex constraints), so $Q$ is the global
        maximum, $f=4$.
  \item $\mathcal{L}=\ln(1+x)+\ln(1+y)+z+\lambda(4-x-y-z)$, folded
        conditions $\mathcal{L}_{x}=\tfrac1{1+x}-\lambda\le0$,
        $\mathcal{L}_y=\tfrac1{1+y}-\lambda\le0$, $\mathcal{L}_z=1-\lambda\le0$,
        each with equality when the variable is positive.
        $(1,1,2)$: budget binds, all variables positive, so $\lambda=\tfrac12$
        (from $x$) and $\lambda=1$ (from $z$): contradiction, not KKT.
        $(0,0,4)$: $z>0$ gives $\lambda=1$; then
        $\mathcal{L}_x=\mathcal{L}_y=0\le0$; budget binds, $\lambda\ge0$:
        \textbf{KKT point}. $(2,2,0)$: $x>0$ gives $\lambda=\tfrac13$, but
        $\mathcal{L}_z=1-\tfrac13>0$: not KKT. $(1,1,1)$: budget slack, so
        $\lambda=0$, but $\mathcal{L}_x=\tfrac12\neq0$: not KKT. The
        objective is concave and the constraints linear: a convex
        programme, so $(0,0,4)$ is the \textbf{global maximum}, $f=4$
        (directly: $\ln(1+t)\le t$ gives $f\le x+y+z\le4$).
  \item $\mathcal{L}=x^{2}+y^{2}+\lambda(4-x-2y)$, folded:
        $\mathcal{L}_x=2x-\lambda\le0$, $\mathcal{L}_y=2y-2\lambda\le0$.
        $(4,0)$: $x>0$ gives $\lambda=8$; $\mathcal{L}_y=-16\le0$: KKT,
        $f=16$. $(0,2)$: $y>0$ gives $\lambda=2$; $\mathcal{L}_x=-2\le0$:
        KKT, $f=4$. $(0,0)$: budget slack, $\lambda=0$,
        $\mathcal{L}_x=\mathcal{L}_y=0$: KKT, $f=0$. $(0.8,1.6)$: budget
        binds, $1.6=\lambda$ and $3.2=2\lambda$: KKT with $\lambda=1.6$,
        $f=3.2$. The maximum of a convex function over the triangle sits at
        a vertex: $f(4,0)=16$ is \textbf{optimal}. Linear constraints make
        KKT necessary, but the programme is not convex (a convex objective
        is maximised), so KKT is not sufficient: here it lets through the
        global minimum $(0,0)$, a boundary minimum and a weaker vertex.
  \item (a) $\mathcal{L}=xy-\tfrac12y^{2}+\mu(6-x-y)+\lambda(y-3)$ with
        $\mu$ free, $\lambda\ge0$. For $x,y>0$: $y-\mu=0$ and
        $x-y-\mu+\lambda=0$. If $\lambda=0$: $\mu=y$, $x=2y$, $3y=6$, $y=2<3$,
        infeasible. So $y=3$, $x=3$, $\mu=3$, and $3-3-3+\lambda=0$ gives
        $\lambda=3\ge0$: the KKT point $(3,3)$, $f=4.5$. It is global: on
        the line $f=(6-y)y-\tfrac12y^{2}=6y-\tfrac32y^{2}$, decreasing for
        $y>2$, so on $3\le y\le6$ the best is $y=3$ (the objective itself is
        not concave, its Hessian has determinant $-1$).
        (b) With $x+y\le6$ ($\mu^{+}\ge0$) and $-x-y\le-6$ ($\mu^{-}\ge0$),
        stationarity contains $\mu^{+}-\mu^{-}$ where (a) had $\mu$, so
        $\mu^{+}-\mu^{-}=3$: every pair $(3+t,t)$, $t\ge0$, works. The
        gradients $(1,1)$ and $(-1,-1)$ are dependent, so LICQ fails and the
        multipliers are not unique; their difference is the free $\mu$ of
        (a). (Non-negativity written as $-x\le0$, $-y\le0$ adds multipliers
        $\nu_x=\nu_y=0$, since both variables are positive.)
        (c) $\lambda=3$ is the cost of the floor: with $y\ge b$ the value is
        $V(b)=6b-\tfrac32b^{2}$, and $V'(3)=-3=-\lambda$: raising the floor by
        $0.1$ costs about $0.3$.
  \item On the plane the unconstrained minimiser is $(1,1,1)$, which
        violates $x\ge2$, so $x\ge2$ binds. $\mathcal{L}=x^{2}+y^{2}+z^{2}
        +\mu(3-x-y-z)+\lambda(2-x)$ (min, $\ge$): $2x=\mu+\lambda$,
        $2y=\mu$, $2z=\mu$. With $x=2$: $y=z$, $2+2y=3$, $y=z=\tfrac12$,
        $\mu=1$, $\lambda=4-1=3\ge0$. Minimum $4+\tfrac14+\tfrac14=4.5$ at
        $(2,\tfrac12,\tfrac12)$, global (convex objective, linear
        constraints).
  \item Without the bound, $x=5$, $y=2.5$ (half the budget on each), which
        violates $x\ge6$. So $x=6$, $y=2$. $\mathcal{L}=\ln x+\ln y
        +\mu(10-x-2y)+\lambda(x-6)$: $\tfrac1y=2\mu$ gives $\mu=\tfrac14$;
        $\tfrac1x-\mu+\lambda=0$ gives $\lambda=\tfrac14-\tfrac16=\tfrac1{12}\ge0$.
        Maximum $\ln12\approx2.485$; global (concave objective, linear
        constraints).
  \item Without $y\le1.5$: $x=y=2$ (closest point of the line to $(1,1)$),
        which violates the bound. So $y=1.5$, $x=2.5$.
        $\mathcal{L}=f+\mu(4-x-y)+\lambda(1.5-y)$: $-2(x-1)-\mu=0$ gives
        $\mu=-3$; $-2(y-1)-\mu-\lambda=0$ gives $-1+3-\lambda=0$, $\lambda=2\ge0$.
        $f=-2.25-0.25=-2.5$. The equality multiplier has no sign
        restriction: $\mu=\mathrm{d}V/\mathrm{d}c$ for the total $c=4$, and
        raising $c$ pushes the bundle further from the bliss point $(1,1)$,
        so the value falls. Indeed $V(c)=-(c-2.5)^{2}-0.25$, $V'(4)=-3$.
  \item \emph{Min.} $\mathcal{L}=x^{2}+2y^{2}+\lambda(3-x-y)$; conditions:
        $\mathcal{L}_x=2x-\lambda\ge0$, $x\ge0$, $x\mathcal{L}_x=0$;
        $\mathcal{L}_y=4y-\lambda\ge0$, $y\ge0$, $y\mathcal{L}_y=0$;
        $x+y\ge3$, $\lambda\ge0$, $\lambda(3-x-y)=0$. $\lambda=0$ would force
        $x=y=0$, infeasible, so $x+y=3$. With $x,y>0$: $x=\tfrac\lambda2$,
        $y=\tfrac\lambda4$, $\tfrac{3\lambda}4=3$: $\lambda=4$, $(2,1)$,
        $f=6$, global (convex programme).
        \emph{Max.} $\mathcal{L}=x^{2}+2y^{2}+\lambda(x+y-3)$; conditions
        $\mathcal{L}_x=2x+\lambda\le0$, $\mathcal{L}_y=4y+\lambda\le0$ with
        $x,y,\lambda\ge0$. Both sums of non-negative terms must be $\le0$,
        so $x=y=\lambda=0$, which violates $x+y\ge3$: \textbf{no KKT point}.
        The constraints are linear, hence the programme is regular and any
        maximum would be a KKT point: so \textbf{no maximum exists} (indeed
        $f\to\infty$ along $x+y\ge3$).
  \item The set is the triangle with vertices $(0,3)$, $(3,3)$, $(3,1.5)$
        ($x\ge0$ is implied).
        \emph{Min.} $\mathcal{L}=f+\lambda(6-x-2y)+\mu_1(x-3)+\mu_2(y-3)$,
        all multipliers $\ge0$ with complementary slackness, stationarity
        $2(x-1)-\lambda+\mu_1=0$, $2(y-2)-2\lambda+\mu_2=0$. $(1,2)$ violates
        $x+2y\ge6$, so try $x+2y=6$ alone: $x-1=\tfrac\lambda2$,
        $y-2=\lambda$, $1+\tfrac\lambda2+4+2\lambda=6$, $\lambda=0.4$:
        $(1.2,2.4)$, inside the other bounds, $f=0.2$. Convex programme:
        global.
        \emph{Max.} $\mathcal{L}=f+\lambda(x+2y-6)+\mu_1(3-x)+\mu_2(3-y)$,
        stationarity $2(x-1)+\lambda-\mu_1=0$, $2(y-2)+2\lambda-\mu_2=0$. The
        largest value among the vertices is $f(3,3)=5$ (against $2$ and
        $4.25$). There $x+2y=9>6$, so $\lambda=0$, and $\mu_1=4$, $\mu_2=2$,
        both $\ge0$: the KKT system holds at the optimal corner.
  \item (a) \textbf{Convex}: convex objective, convex $g=x^{2}+y^{2}$ in a
        $\le$ constraint. (b) \textbf{Not}: a $\ge$ constraint needs a
        concave function, $x^{2}+y^{2}$ is convex; the set outside the disc
        is not convex. (c) \textbf{Not}: $x^{2}-y^{2}$ is not concave; the
        set consists of two separate pieces. (d) \textbf{Not}: $y-x^{2}\le0$
        has a concave, not convex, left side; the region under a convex
        parabola is not convex (the objective is fine). (e)
        \textbf{Convex}: the objective has Hessian
        $\begin{pmatrix}12x^{2}+2&-2\\-2&2\end{pmatrix}$, PSD everywhere
        (Additional Problem~38), and the constraint is linear. (f)
        \textbf{Not}: maximising requires a concave objective; this is
        the disc of (a) with a convex objective maximised.
  \item On the feasible set $x^{2}\ge1+y^{2}$, so
        $f=x^{2}+y^{2}\ge1+2y^{2}\ge1$, with equality iff $y=0$, $x^{2}=1$:
        \textbf{two minimisers} $(\pm1,0)$, $f=1$. KKT check:
        $\mathcal{L}=x^{2}+y^{2}+\lambda(1-x^{2}+y^{2})$: $2x(1-\lambda)=0$,
        $2y(1+\lambda)=0$. So $y=0$ ($\lambda\ge0$), and $x\neq0$ forces
        $\lambda=1$: both points with $\lambda=1$. Two optima are possible
        because the programme is not convex.
  \item Convexity: $f_{xx}=12x^{2}+2>0$, $f_{yy}=2>0$,
        $\det=24x^{2}+4-4=24x^{2}\ge0$: PSD everywhere. Unconstrained,
        $\grad f=(4x^{3}+2x-2y,\ -2x+2y)=0$ gives $y=x$, $4x^{3}=0$: the
        origin, which violates $x+y\ge3$. So the constraint binds:
        $\mathcal{L}=f+\lambda(3-x-y)$, $4x^{3}+2x-2y=\lambda$,
        $-2x+2y=\lambda$, $y=3-x$. From the second, $\lambda=6-4x$; the
        first becomes $4x^{3}+2x-6+2x=6-4x$, i.e.\ $x^{3}+2x-3=0
        =(x-1)(x^{2}+x+3)$, so $x=1$, $y=2$, $\lambda=2\ge0$. Minimum
        $f(1,2)=1+1-4+4=2$, global (convex programme).
  \item Min form: stationarity $2(x-2)+\mu_1+\mu_2=0$,
        $2(y-2)+\mu_1+2\mu_2=0$, $2(z-2)+\mu_1=0$. Cases: none active,
        $(2,2,2)$: sum $6>3$, infeasible. Only $g_1$: $(1,1,1)$ violates
        $x+2y\le2$. Only $g_2$: $(\tfrac65,\tfrac25,2)$, sum $3.6>3$,
        infeasible. Both: $x=2-\tfrac{\mu_1+\mu_2}2$,
        $y=2-\tfrac{\mu_1+2\mu_2}2$, $z=2-\tfrac{\mu_1}2$; the constraints
        give $\mu_1+\mu_2=2$ and $\mu_1+2\mu_2=3$, so $\mu_1=\mu_2=1$ and
        $(x,y,z)=(1,\tfrac12,\tfrac32)$, $f=1+\tfrac94+\tfrac14=\tfrac72$.
        Multipliers positive, convex programme: \textbf{global minimum}.
  \item $\grad f=(y+z,\,x+z,\,x+y)=H\mathbf{v}$ with
        $H=\begin{pmatrix}0&1&1\\1&0&1\\1&1&0\end{pmatrix}$, which has
        eigenvalues $2$ (eigenvector $(1,1,1)$) and $-1$ (twice): $f$ is
        \textbf{not concave}. KKT: $H\mathbf{v}=2\lambda\mathbf{v}$,
        $\lambda\ge0$, $\lambda(3-\|\mathbf{v}\|^{2})=0$. $\lambda=0$: $H$ is
        invertible, so $\mathbf{v}=0$, $f=0$. Ball boundary: $\mathbf{v}$ is
        an eigenvector with $2\lambda\in\{2,-1\}$; only $\lambda=1$ is
        allowed, with $\mathbf{v}\parallel(1,1,1)$, $\|\mathbf{v}\|^{2}=3$:
        $\pm(1,1,1)$, $f=3$. Symmetry explains the shortcut: $f$ is
        unchanged by permuting coordinates and by $\mathbf{v}\mapsto-\mathbf{v}$,
        so try $x=y=z$ and expect solutions in pairs. The ball is compact
        and the programme regular (Slater: the origin), so the maximum is
        among the KKT points: $f=3$ at the \textbf{two maximisers}
        $(1,1,1)$ and $(-1,-1,-1)$.
  \item $U$ is increasing, so the budget binds:
        $\mathcal{L}=U+\lambda(1.5-x-y-z)+\mu(x-1)$, with
        $\mathcal{L}_x=1-\lambda+\mu=0$, $\mathcal{L}_y=\tfrac1{\sqrt y}-\lambda$
        ($=0$ if $y>0$), $\mathcal{L}_z=\tfrac1{1+z}-\lambda\le0$ ($=0$ if
        $z>0$). Case $x>1$: $\mu=0$, $\lambda=1$, $y=1$, and $z=0$; then
        $x=0.5<1$: infeasible. So $x=1$ and $y+z=0.5$. If $z>0$:
        $\sqrt y=1+z$, so $(1+z)^{2}+z=0.5$, $z^{2}+3z+0.5=0$, both roots
        negative: impossible. Hence $z=0$, $y=0.5$, $\lambda=\sqrt2$;
        $\mathcal{L}_z=1-\sqrt2\le0$ \checkmark, $\mu=\sqrt2-1\approx0.414\ge0$
        \checkmark. Optimum $(1,\tfrac12,0)$, $U=1+\sqrt2\approx2.414$,
        global (concave $U$, linear constraints).
  \item (a) $\max x^{1/2}y^{1/2}$ s.t.\ $2x+y\le40$ (money), $x+2y\le40$
        (coupons), $x,y\ge0$. (b) Money alone: half on each good,
        $x=10$, $y=20$, needing $10+40=50>40$ coupons. Coupons alone: half
        of the coupons on each, $x=20$, $y=10$, costing $40+10=50>40$ euros.
        (c) Both binding: $2x+y=40$, $x+2y=40$ give $x=y=\tfrac{40}3$. There
        $U_x=U_y=\tfrac12$, and
        $(\tfrac12,\tfrac12)=\lambda_m(2,1)+\lambda_c(1,2)$ gives
        $\lambda_m=\lambda_c=\tfrac16>0$: KKT, and global (concave $U$,
        linear constraints). An extra euro and an extra coupon are each
        worth $\tfrac16$ util, so one coupon is worth
        $\lambda_c/\lambda_m=1$ euro: paying $1.50$ is \textbf{not}
        worthwhile.
  \item Unconstrained: $x=10$, $y=8$, needing $18>10$ units of $A$. Only
        $A$: $20-2x=16-2y$ and $x+y=10$ give $(6,4)$, but $2\cdot6+4=16>15$.
        Only $B$: $20-2x=2\lambda$, $16-2y=\lambda$, so $x=2y-6$, and
        $2x+y=15$ gives $(4.8,5.4)$, but $4.8+5.4=10.2>10$. Both binding:
        $x+y=10$, $2x+y=15$, so $(5,5)$, with $\grad\pi=(10,6)
        =\lambda_A(1,1)+\lambda_B(2,1)$: $\lambda_B=4$, $\lambda_A=2$, both
        positive. Profit $100-25+80-25=130$, global (concave, linear
        constraints). One more unit of $A$ raises profit by about $2$, one
        more unit of $B$ by about $4$: $B$ is the scarcer resource.
  \item Unconstrained, $3\mathrm{e}^{-a_1/20}=1$ and $4\mathrm{e}^{-a_2/10}=1$:
        $a_1=20\ln3\approx21.97$, $a_2=10\ln4\approx13.86$, total
        $\approx35.84$. \emph{$B=50$:} the budget is slack, $\lambda=0$, and
        $\pi=40+30-35.84\approx34.16$; $14.16$ of the budget is left
        unspent, because beyond that point an extra euro of advertising
        earns less than a euro. \emph{$B=30$:} the budget binds,
        $3\mathrm{e}^{-a_1/20}-1=\lambda=4\mathrm{e}^{-a_2/10}-1$, so
        $a_1=20\ln\tfrac{3}{1+\lambda}$, $a_2=10\ln\tfrac{4}{1+\lambda}$, and
        $a_1+a_2=30$ gives $\ln(1+\lambda)=\tfrac{20\ln3+10\ln4-30}{30}$,
        $\lambda\approx0.215$, $a_1\approx18.08$, $a_2\approx11.92$,
        $\pi\approx33.56$. Now an extra euro of budget would raise profit by
        about $0.21$. ($\pi$ is concave, so both answers are global.)
  \item On the line $f=x(10-x)$ for $x\in[2,6]$. $\mathcal{L}$ for the max:
        $xy+\mu(10-x-y)+\lambda_1(x-2)+\lambda_2(6-x)$: at $(5,5)$ both
        bounds slack, $\mu=5$, \textbf{maximum} $25$. For the min,
        $\mathcal{L}=xy+\mu(10-x-y)+\lambda_1(2-x)+\lambda_2(x-6)$,
        stationarity $y-\mu-\lambda_1+\lambda_2=0$, $x-\mu=0$. At $(2,8)$:
        $\mu=2$, $\lambda_1=6\ge0$ (the lower bound binds), $f=16$. At
        $(6,4)$: $\mu=6$, $\lambda_2=2\ge0$ (the upper bound binds), $f=24$:
        also a KKT point, a local minimum only. \textbf{Minimum} $16$ at
        $(2,8)$; the lower bound binds at the minimum, no bound at the
        maximum.
  \item Substitute $x=u-v$ ($u,v\ge0$, $x$ free) and $z=-w$ ($w\ge0$).
        (a) $\max\,(u-v)y-w^{2}-2w$ s.t.\ $5-(u-v)^{2}-y\ge0$,
        $(u-v)-y+1\ge0$, $u,v,y,w\ge0$.
        (b) $\min\,-(u-v)y+w^{2}+2w$ s.t.\ $(u-v)^{2}+y-5\le0$,
        $-(u-v)+y-1\le0$, $u,v,y,w\ge0$.
        (In the chapter's own form, $\max f$ s.t.\ $g_i\le b_i$, keep
        $x$ free and write $-y\le0$, $z\le0$, $-x+y\le1$.)
  \item Substitute $x=u-v$, $w=-s$ ($u,v,s\ge0$) and split the equality.
        Form ``$\max$, $\ge0$'': $\max\,-(u-v-1)^{2}-\mathrm{e}^{-s}$ s.t.\
        $u-v-s-2\ge0$, $-u+v+s+2\ge0$, $u,v,s\ge0$. Form ``$\min$, $\le0$'':
        $\min\,(u-v-1)^{2}+\mathrm{e}^{-s}$ s.t.\ $u-v-s-2\le0$,
        $-u+v+s+2\le0$, $u,v,s\ge0$. Solution: $x=2-w$ gives
        $\varphi(w)=(1-w)^{2}+\mathrm{e}^{w}$ on $w\le0$, with
        $\varphi'(w)=-2(1-w)+\mathrm{e}^{w}\le-2+1<0$ there: decreasing, so
        the bound binds, $w^*=0$, $x^*=2$, minimum $2$.
  \item (a) Side $\tfrac x4$, radius $\tfrac{20-x}{2\pi}$:
        $A(x)=\tfrac{x^{2}}{16}+\tfrac{(20-x)^{2}}{4\pi}$ on $[0,20]$, with
        $A''=\tfrac18+\tfrac1{2\pi}>0$, convex. $A'(x)=\tfrac x8-\tfrac{20-x}{2\pi}=0$
        gives $x=\tfrac{80}{4+\pi}\approx11.20$, the \textbf{minimum}
        $A=\tfrac{100}{4+\pi}\approx14.00$. The maximum of a convex function
        sits at an endpoint: $A(0)=\tfrac{100}{\pi}\approx31.83$ (all circle)
        beats $A(20)=25$, so the \textbf{maximum} is at $x=0$.
        (b) $\mathcal{L}=\tfrac{x^{2}}{16}+\tfrac{y^{2}}{4\pi}+\lambda(20-x-y)$:
        $\tfrac x8=\lambda=\tfrac{y}{2\pi}$, so $y=\tfrac{\pi}{4}x$ and the
        same $x=\tfrac{80}{4+\pi}$. $\det\overline{H}=-\mathcal{L}_{yy}-\mathcal{L}_{xx}
        =-\tfrac1{2\pi}-\tfrac18<0$: a \textbf{minimum}, the only Lagrange
        point. (c) On $\{x+y=20,\ x,y\ge0\}$ the corners carry
        bound multipliers: with
        $\mathcal{L}=A+\mu(20-x-y)+\nu_xx+\nu_yy$, at $(0,20)$:
        $\mu=\tfrac{20}{2\pi}=\tfrac{10}\pi$ and $\nu_x=\mu>0$; at $(20,0)$:
        $\mu=\tfrac52=\nu_y$. Both corners are KKT points of the maximisation,
        and the better one, $(0,20)$, is the maximum. Lagrange looks only for
        tangencies on the line and never sees the bounds $x,y\ge0$.
  \item Max: $\mathcal{L}=x^{2}+xy+y^{2}+\lambda(x+y-2)$, folded
        $\mathcal{L}_x=2x+y+\lambda\le0$, $\mathcal{L}_y=x+2y+\lambda\le0$
        with $x,y,\lambda\ge0$: only $x=y=\lambda=0$, which is infeasible.
        No KKT point; the constraints are linear, so the programme is
        regular and \textbf{no maximum exists} ($f\to\infty$).
        Min: $\mathcal{L}=f+\lambda(2-x-y)$, $2x+y=\lambda=x+2y$ for
        $x,y>0$, so $x=y=1$, $\lambda=3\ge0$, $f=3$; convex objective
        ($\Hess=\begin{pmatrix}2&1\\1&2\end{pmatrix}$ PD), global.
  \item (a) Objective $f(p)=\|p-a\|^{2}+\|p-b\|^{2}$, total squared
        distance to the shops; the constraint function $\|p\|^{2}$ measures
        the squared distance from the centre of the zone, and
        $\|p\|^{2}\le25$ says the warehouse lies in the zone. (b) A disc is
        convex; $f$ is a convex quadratic and $\|p\|^{2}$ convex: a
        \textbf{convex programme}; Slater holds at the origin:
        \textbf{regular}. (c) Regular, so KKT is \textbf{necessary}; convex,
        so it is \textbf{sufficient}. A Lagrange point on the circle with
        $\lambda\ge0$ is therefore the global optimum, and it is unique
        because $f$ is strictly convex. (d) $f=2\|p-m\|^{2}+8$ with midpoint
        $m=(6,2)$, $\|m\|=\sqrt{40}>5$: the unconstrained optimum lies
        outside, so the circle binds, and the optimum is the point of the
        disc nearest to $m$: $p^*=\tfrac{5}{\sqrt{40}}m\approx(4.743,\,1.581)$,
        $f^*=2(\sqrt{40}-5)^{2}+8\approx11.51$, with
        $\lambda=0.4(\sqrt{40}-5)\approx0.53\ge0$.
  \item (a) Linear constraint: necessary; concave objective: a convex
        programme, so also \textbf{sufficient}. (b) Linear constraints:
        \textbf{necessary}; a convex objective is maximised, so
        \textbf{not sufficient} ($(0,0)$ is a KKT point and the minimum).
        (c) Not regular at the cusp $(1,0)$, and that is the optimum:
        \textbf{not necessary} (no point at all satisfies KKT here); the
        feasible set is not convex, so no sufficiency either. (d)
        $\grad(x^{2}+y^{2})\neq0$ on the circle: regular, \textbf{necessary};
        a $\ge$ constraint with a convex function makes the programme
        non-convex, so sufficiency is \textbf{not guaranteed} by the theorem
        (here every KKT point, the whole unit circle, happens to be optimal).
  \item If $\grad f=\sum\lambda_i\grad g_i=\sum\lambda_i'\grad g_i$ over the
        active constraints, then $\sum(\lambda_i-\lambda_i')\grad g_i=0$, and
        linear independence forces $\lambda_i=\lambda_i'$: at a regular point
        the multipliers are unique, and each is the shadow price of its own
        constraint. Non-unique multipliers signal dependent active
        gradients: a corner with more active constraints than dimensions
        (the origin in Additional Problem~25, three constraints in the
        plane), or an equality written as two inequalities (Additional
        Problem~30(b)). Then a single $\lambda_i$ no longer measures the
        value of relaxing constraint $i$ alone.
  \item (a) $V(x_1,x_2)=100x_1^{2}+36x_2^{2}+16(1-x_1-x_2)^{2}$, return
        $8x_1+5x_2+3(1-x_1-x_2)\ge6$, i.e.\ $5x_1+2x_2\ge3$, and $x_1,x_2\ge0$,
        $x_1+x_2\le1$. $\mathcal{L}=V+\lambda(3-5x_1-2x_2)-\nu_1x_1-\nu_2x_2+\nu_3(x_1+x_2-1)$,
        all multipliers $\ge0$ with complementary slackness, and
        $V_{x_1}-5\lambda-\nu_1+\nu_3=0$, $V_{x_2}-2\lambda-\nu_2+\nu_3=0$,
        where $V_{x_1}=200x_1-32x_3$, $V_{x_2}=72x_2-32x_3$.
        (b) With $\nu_1=\nu_2=\nu_3=0$: $200x_1-32x_3=5\lambda$,
        $72x_2-32x_3=2\lambda$, $5x_1+2x_2=3$. Solving,
        $x_1=\tfrac{147}{361}\approx0.407$, $x_2=\tfrac{174}{361}\approx0.482$,
        $x_3=\tfrac{40}{361}\approx0.111$, $\lambda=\tfrac{296}{19}\approx15.58\ge0$.
        Checks: $x_1,x_2>0$ and $x_1+x_2<1$, so the slack constraints are
        indeed slack, as their zero multipliers require. Minimal variance
        $\tfrac{9076}{361}\approx25.14$, $\sigma\approx5.01\%$.
        (c) $\Hess V=\begin{pmatrix}232&32\\32&104\end{pmatrix}$ is PD, the
        constraints linear: a convex programme, so the KKT point is the
        unique global minimum. $\lambda\approx15.6$ is the extra variance per
        extra percentage point of required return.
  \item With target $7.5\%$, i.e.\ $5x_1+2x_2\ge4.5$, the same guess gives
        $x_1\approx0.644$, $x_2\approx0.640$, $x_3\approx-0.284<0$: short
        selling, infeasible. So also $x_1+x_2\le1$ binds ($x_3=0$):
        $5x_1+2x_2=4.5$, $x_1+x_2=1$ give $x_1=\tfrac56$, $x_2=\tfrac16$.
        Stationarity with $x_3=0$: $200\cdot\tfrac56-5\lambda+\nu_3=0$,
        $72\cdot\tfrac16-2\lambda+\nu_3=0$; subtracting,
        $3\lambda=\tfrac{1000}{6}-12$, $\lambda=\tfrac{464}9\approx51.6$,
        $\nu_3=2\lambda-12=\tfrac{820}9\approx91.1$, both $\ge0$; $\nu_1=\nu_2=0$
        with $x_1,x_2>0$. All conditions hold; variance
        $\tfrac{2500}{36}+1=\tfrac{634}9\approx70.4$. Global by the same
        convexity argument.
  \item (a) $\pi_{x_1}=\tfrac{6x_2}{1+x_1x_2}-1$,
        $\pi_{x_2}=\tfrac{6x_1}{1+x_1x_2}-1$: both zero gives $x_1=x_2=t$
        with $6t=1+t^{2}$, $t=3\pm2\sqrt2$. At $(t,t)$,
        $\Hess\pi=\tfrac{6}{(1+t^{2})^{2}}\begin{pmatrix}-t^{2}&1\\1&-t^{2}\end{pmatrix}$,
        with determinant of the sign of $t^{4}-1$: $t_0=3-2\sqrt2\approx0.172$
        gives a \textbf{saddle} ($\pi\approx-0.169$), $t_1=3+2\sqrt2\approx5.83$
        a \textbf{local maximum} ($\pi\approx9.670$).
        (b) $B=20$: $(t_0,t_0)$ is feasible with slack budget; $\lambda=0$
        and $\grad\pi=0$ satisfy every KKT condition, but it is a saddle. The
        candidates: $(t_1,t_1)$ with $2t_1\approx11.66<20$, $\pi\approx9.67$;
        the origin, $\pi=0$; on the budget line the best point is
        $(10,10)$ (symmetry), $\pi=6\ln101-20\approx7.69$, and its multiplier
        $\tfrac{12B}{4+B^{2}}-1<0$ rules it out anyway; on the axes
        $\pi=-x_i\le0$. Optimum $(t_1,t_1)$, budget slack, $\lambda=0$:
        $8.34$ of the budget stays unspent.
        (c) $B=4<2t_1$: on the budget line $x_1x_2$ is largest at $(2,2)$,
        $\pi=6\ln5-4\approx5.66$, with
        $\lambda=\tfrac{6x_2}{1+x_1x_2}-1=\tfrac{12}{5}-1=1.4>0$: an extra
        euro of budget is worth $1.40$. This beats the origin and the
        saddle, so it is the optimum. $B=0.5$: the budget point
        $(0.25,0.25)$ is a KKT point ($\lambda\approx0.41$) but
        $\pi=6\ln1.0625-0.5\approx-0.14<0$; the origin is a KKT point too
        (budget slack, $\lambda=0$, $\pi_{x_i}=-1\le0$) with $\pi=0$:
        \textbf{shut down}. The switch happens where
        $6\ln(1+B^{2}/4)=B$, at $B\approx0.71$. A non-convex programme: several
        KKT points, and the regime depends on $B$.
  \item $R$ increases in $x_1$ whenever $x_2>0$ (and vice versa), and
        $R=0$ on the axes, so any unspent budget can raise revenue: the
        budget binds and may be written as $x_1+x_2=B$. Then $x_1x_2$ is
        maximal at $x_1=x_2=\tfrac B2$, $R=6\ln(1+\tfrac{B^{2}}4)$ and
        $\lambda=R_{x_1}=\tfrac{12B}{4+B^{2}}$. $B=4$: $(2,2)$,
        $R\approx9.66$, $\lambda=2.4$. $B=20$: $(10,10)$, $R=6\ln101\approx27.69$,
        $\lambda\approx0.594$. Unlike the profit problem, where input costs
        enter the objective and the budget is left partly unspent for large
        $B$ (and $\lambda=0$), revenue maximisation always exhausts the
        budget with $\lambda>0$.
  \item Domain constraint $g=x^{2}+y^{2}\le5$; KKT for the maximum:
        $\grad f=\lambda\grad g$, $\lambda\ge0$, $\lambda(5-g)=0$. Try the
        domain constraint slack, $\lambda=0$: with $R=\sqrt{5-x^{2}-y^{2}}$,
        $f_x=1-\tfrac xR=0$, $f_y=2-\tfrac yR=0$, so $x=R$, $y=2R$ and
        $R^{2}=5-5R^{2}$, $R=\sqrt{5/6}$: $(x,y)\approx(0.913,\,1.826)$ with
        $x^{2}+y^{2}=\tfrac{25}6<5$, slack indeed. All KKT conditions hold
        with $\lambda=0$, and $f^*=6R=\sqrt{30}\approx5.477$. Global: $x+2y$
        is linear and $\sqrt{5-x^{2}-y^{2}}$ is concave (an increasing
        concave function of a concave function), so $f$ is concave on the
        convex domain. Near the boundary the slope of the square root becomes unbounded,
        which keeps the optimum away from it.
  \item Constraint $g=x^{2}-y\le0$ with $\grad g=(2x,-1)\neq0$: LICQ holds
        everywhere, the programme is regular. KKT for the maximum:
        $\grad f=\lambda\grad g$,
        \[
          \Bigl(1-3x\sqrt{y-x^{2}},\ \ 1+\tfrac32\sqrt{y-x^{2}}\Bigr)
          =\lambda\,(2x,\,-1).
        \]
        The second components give $\lambda=-1-\tfrac32\sqrt{y-x^{2}}\le-1<0$,
        violating $\lambda\ge0$ (and $\lambda=0$, the interior case, is
        excluded too). No point satisfies the KKT conditions; by regularity
        a maximum would have to, so \textbf{no maximum exists}. Directly:
        $f(0,y)=y+y^{3/2}\to\infty$.
\end{solutions}
